% !TeX root = ../main.tex
\clearpage
\section*{附录3\quad 核心代码（Python节选）}
\phantomsection\label{appendix:code}
\addcontentsline{toc}{section}{附录3：核心代码（Python节选）}

以下代码按“读入—建模—求解—复核”的顺序排列，与正文两段伪代码相对应。各段取自当前工程的连续源码，保留原行号，排版仅去除共同的前导缩进；原样节选、完整文件、依赖模块和指纹清单随LaTeX源包提供。代码需要原工程的配置、输入数据和调用环境，不能单独运行。当前源码用于说明实现，三城历史结果仍以其原始实验记录为准。

\begin{table}[H]
\centering\small
\caption{附录源码与正文方法的对应关系}
\begin{tabularx}{\textwidth}{@{}cL{0.31\textwidth}Y@{}}
\toprule
序号 & 文件与原行号 & 对应内容\\\midrule
1 & \texttt{unified\_cli.py}，93--118 & 配置读入、构模、求解与导出\\
2 & \texttt{building\_thermal.py}，31--46 & 一阶建筑热状态离散\\
3 & \texttt{model.py}，1603--1697 & 节点—温位—小时热平衡\\
4 & \texttt{model.py}，1425--1471 & 储热状态与边界\\
5 & \texttt{linear.py}，177--205 & 稀疏MILP求解接口\\
6 & \texttt{experiments.py}，43--65 & 热源与储热投资冻结节选\\
7 & \texttt{heating\_season.py}，183--246 & 连续供暖季复算及检查\\
8 & \texttt{experiments.py}，351--411 & 独立碳帽扫描与样本筛选\\
\bottomrule
\end{tabularx}
\end{table}

正文算法\ref{alg:season-workflow}概括主流程；其中两阶段接口由第6、7段衔接。算法\ref{alg:epsilon-workflow}对应第8段，单独清除原碳帽并将碳价置零。该接口与保留碳价和偏好的历史赛事扫描分别登记。

\Needspace{14\baselineskip}
\subsection*{1．入口与导出}
读取配置及时间序列，构造模型并导出结果。这是主入口的单次求解分支，其他分支见完整文件。

\begingroup
\setstretch{1.0}
\begin{Verbatim}[commandchars=\\\{\},numbers=left,firstnumber=93,stepnumber=1,fontsize=\fontsize{9.5}{11.5}\selectfont,breaklines=true,breakanywhere=true,frame=single,framesep=3mm,numbersep=5pt,rulecolor=\color{black},xleftmargin=18pt]
\PYG{k}{else}\PYG{p}{:}
    \PYG{n}{config}\PYG{p}{,} \PYG{n}{frame} \PYG{o}{=} \PYG{n}{load\PYGZus{}inputs}\PYG{p}{(}\PYG{n}{config\PYGZus{}path}\PYG{p}{,} \PYG{n}{timeseries\PYGZus{}path}\PYG{p}{)}
    \PYG{k}{if} \PYG{n}{scenario\PYGZus{}specs\PYGZus{}path}\PYG{p}{:}
        \PYG{n}{specs} \PYG{o}{=} \PYG{n}{yaml}\PYG{o}{.}\PYG{n}{safe\PYGZus{}load}\PYG{p}{(}\PYG{n}{scenario\PYGZus{}specs\PYGZus{}path}\PYG{o}{.}\PYG{n}{read\PYGZus{}text}\PYG{p}{(}\PYG{n}{encoding}\PYG{o}{=}\PYG{l+s+s2}{\PYGZdq{}}\PYG{l+s+s2}{utf\PYGZhy{}8}\PYG{l+s+s2}{\PYGZdq{}}\PYG{p}{)}\PYG{p}{)}
        \PYG{n}{demand\PYGZus{}columns} \PYG{o}{=} \PYG{p}{[}
            \PYG{n}{column} \PYG{k}{for} \PYG{n}{node} \PYG{o+ow}{in} \PYG{n}{config}\PYG{p}{[}\PYG{l+s+s2}{\PYGZdq{}}\PYG{l+s+s2}{nodes}\PYG{l+s+s2}{\PYGZdq{}}\PYG{p}{]}
            \PYG{k}{for} \PYG{n}{column} \PYG{o+ow}{in} \PYG{n}{node}\PYG{o}{.}\PYG{n}{get}\PYG{p}{(}\PYG{l+s+s2}{\PYGZdq{}}\PYG{l+s+s2}{demand\PYGZus{}columns}\PYG{l+s+s2}{\PYGZdq{}}\PYG{p}{,} \PYG{p}{\PYGZob{}}\PYG{p}{\PYGZcb{}}\PYG{p}{)}\PYG{o}{.}\PYG{n}{values}\PYG{p}{(}\PYG{p}{)}
        \PYG{p}{]}
        \PYG{n}{frame} \PYG{o}{=} \PYG{n}{build\PYGZus{}correlated\PYGZus{}northern\PYGZus{}scenario\PYGZus{}bundle}\PYG{p}{(}
            \PYG{n}{frame}\PYG{p}{,} \PYG{n}{demand\PYGZus{}columns}\PYG{p}{,} \PYG{n}{specs}\PYG{p}{[}\PYG{l+s+s2}{\PYGZdq{}}\PYG{l+s+s2}{scenarios}\PYG{l+s+s2}{\PYGZdq{}}\PYG{p}{]}\PYG{p}{,}
            \PYG{n}{period\PYGZus{}weight}\PYG{o}{=}\PYG{n+nb}{float}\PYG{p}{(}\PYG{n}{specs}\PYG{o}{.}\PYG{n}{get}\PYG{p}{(}\PYG{l+s+s2}{\PYGZdq{}}\PYG{l+s+s2}{period\PYGZus{}weight}\PYG{l+s+s2}{\PYGZdq{}}\PYG{p}{,} \PYG{n}{config}\PYG{p}{[}\PYG{l+s+s2}{\PYGZdq{}}\PYG{l+s+s2}{planning}\PYG{l+s+s2}{\PYGZdq{}}\PYG{p}{]}\PYG{o}{.}\PYG{n}{get}\PYG{p}{(}\PYG{l+s+s2}{\PYGZdq{}}\PYG{l+s+s2}{representative\PYGZus{}day\PYGZus{}weight}\PYG{l+s+s2}{\PYGZdq{}}\PYG{p}{,} \PYG{l+m+mf}{1.0}\PYG{p}{)}\PYG{p}{)}\PYG{p}{)}\PYG{p}{,}
        \PYG{p}{)}
        \PYG{n}{config} \PYG{o}{=} \PYG{n}{copy}\PYG{o}{.}\PYG{n}{deepcopy}\PYG{p}{(}\PYG{n}{config}\PYG{p}{)}
        \PYG{n}{config}\PYG{p}{[}\PYG{l+s+s2}{\PYGZdq{}}\PYG{l+s+s2}{planning}\PYG{l+s+s2}{\PYGZdq{}}\PYG{p}{]}\PYG{p}{[}\PYG{l+s+s2}{\PYGZdq{}}\PYG{l+s+s2}{risk}\PYG{l+s+s2}{\PYGZdq{}}\PYG{p}{]} \PYG{o}{=} \PYG{n+nb}{dict}\PYG{p}{(}\PYG{n}{specs}\PYG{o}{.}\PYG{n}{get}\PYG{p}{(}\PYG{l+s+s2}{\PYGZdq{}}\PYG{l+s+s2}{risk}\PYG{l+s+s2}{\PYGZdq{}}\PYG{p}{,} \PYG{p}{\PYGZob{}}\PYG{p}{\PYGZcb{}}\PYG{p}{)}\PYG{p}{)}
        \PYG{n}{validate\PYGZus{}inputs}\PYG{p}{(}\PYG{n}{config}\PYG{p}{,} \PYG{n}{frame}\PYG{p}{)}
    \PYG{n}{model} \PYG{o}{=} \PYG{n}{HeatPlanningModel}\PYG{p}{(}\PYG{n}{config}\PYG{p}{,} \PYG{n}{frame}\PYG{p}{)}
    \PYG{n}{artifacts} \PYG{o}{=} \PYG{n}{model}\PYG{o}{.}\PYG{n}{solve}\PYG{p}{(}
        \PYG{n}{mip\PYGZus{}gap}\PYG{o}{=}\PYG{n}{args}\PYG{o}{.}\PYG{n}{mip\PYGZus{}gap}\PYG{p}{,}
        \PYG{n}{time\PYGZus{}limit}\PYG{o}{=}\PYG{n}{args}\PYG{o}{.}\PYG{n}{time\PYGZus{}limit}\PYG{p}{,}
    \PYG{p}{)}
    \PYG{n}{export\PYGZus{}results}\PYG{p}{(}\PYG{n}{profile\PYGZus{}output}\PYG{p}{,} \PYG{n}{artifacts}\PYG{p}{,} \PYG{n}{config}\PYG{p}{,} \PYG{n}{model}\PYG{o}{.}\PYG{n}{ts}\PYG{p}{)}
    \PYG{n}{summary} \PYG{o}{=} \PYG{p}{\PYGZob{}}\PYG{n}{key}\PYG{p}{:} \PYG{n}{value} \PYG{k}{for} \PYG{n}{key}\PYG{p}{,} \PYG{n}{value} \PYG{o+ow}{in} \PYG{n}{artifacts}\PYG{o}{.}\PYG{n}{summary}\PYG{o}{.}\PYG{n}{items}\PYG{p}{(}\PYG{p}{)} \PYG{k}{if} \PYG{o+ow}{not} \PYG{n+nb}{isinstance}\PYG{p}{(}\PYG{n}{value}\PYG{p}{,} \PYG{n+nb}{dict}\PYG{p}{)}\PYG{p}{\PYGZcb{}}
    \PYG{p}{(}\PYG{n}{profile\PYGZus{}output} \PYG{o}{/} \PYG{l+s+s2}{\PYGZdq{}}\PYG{l+s+s2}{scenario\PYGZus{}results.json}\PYG{l+s+s2}{\PYGZdq{}}\PYG{p}{)}\PYG{o}{.}\PYG{n}{write\PYGZus{}text}\PYG{p}{(}
        \PYG{n}{json}\PYG{o}{.}\PYG{n}{dumps}\PYG{p}{(}\PYG{n}{artifacts}\PYG{o}{.}\PYG{n}{summary}\PYG{o}{.}\PYG{n}{get}\PYG{p}{(}\PYG{l+s+s2}{\PYGZdq{}}\PYG{l+s+s2}{scenario\PYGZus{}results}\PYG{l+s+s2}{\PYGZdq{}}\PYG{p}{,} \PYG{p}{\PYGZob{}}\PYG{p}{\PYGZcb{}}\PYG{p}{)}\PYG{p}{,} \PYG{n}{ensure\PYGZus{}ascii}\PYG{o}{=}\PYG{k+kc}{False}\PYG{p}{,} \PYG{n}{indent}\PYG{o}{=}\PYG{l+m+mi}{2}\PYG{p}{)}\PYG{p}{,}
        \PYG{n}{encoding}\PYG{o}{=}\PYG{l+s+s2}{\PYGZdq{}}\PYG{l+s+s2}{utf\PYGZhy{}8}\PYG{l+s+s2}{\PYGZdq{}}\PYG{p}{,}
    \PYG{p}{)}
\end{Verbatim}

\endgroup

\Needspace{14\baselineskip}
\subsection*{2．建筑热状态离散}
由热损失系数、热容和时间步长计算一阶热状态的精确离散系数。该函数服务于独立逐楼热模型。

\begingroup
\setstretch{1.0}
\begin{Verbatim}[commandchars=\\\{\},numbers=left,firstnumber=31,stepnumber=1,fontsize=\fontsize{9.5}{11.5}\selectfont,breaklines=true,breakanywhere=true,frame=single,framesep=3mm,numbersep=5pt,rulecolor=\color{black},xleftmargin=18pt]
\PYG{k}{def}\PYG{+w}{ }\PYG{n+nf}{exact\PYGZus{}rc\PYGZus{}coefficients}\PYG{p}{(}
    \PYG{n}{heat\PYGZus{}loss\PYGZus{}kw\PYGZus{}per\PYGZus{}c}\PYG{p}{:} \PYG{n+nb}{float}\PYG{p}{,} \PYG{n}{thermal\PYGZus{}capacity\PYGZus{}kwh\PYGZus{}per\PYGZus{}c}\PYG{p}{:} \PYG{n+nb}{float}\PYG{p}{,} \PYG{n}{dt\PYGZus{}hours}\PYG{p}{:} \PYG{n+nb}{float}\PYG{p}{,}
\PYG{p}{)} \PYG{o}{\PYGZhy{}}\PYG{o}{\PYGZgt{}} \PYG{n+nb}{tuple}\PYG{p}{[}\PYG{n+nb}{float}\PYG{p}{,} \PYG{n+nb}{float}\PYG{p}{]}\PYG{p}{:}
\PYG{+w}{    }\PYG{l+s+sd}{\PYGZdq{}\PYGZdq{}\PYGZdq{}Return a, b for T\PYGZus{}end = a*T\PYGZus{}start + (1\PYGZhy{}a)*Tout + b*(heat\PYGZhy{}cool+gains).}

\PYG{l+s+sd}{    expm1 avoids cancellation at small H*dt/C. H=0 is deliberately rejected:}
\PYG{l+s+sd}{    this module requires a positive envelope/ventilation heat loss coefficient.}
\PYG{l+s+sd}{    \PYGZdq{}\PYGZdq{}\PYGZdq{}}
    \PYG{n}{h} \PYG{o}{=} \PYG{n}{\PYGZus{}number}\PYG{p}{(}\PYG{n}{heat\PYGZus{}loss\PYGZus{}kw\PYGZus{}per\PYGZus{}c}\PYG{p}{,} \PYG{l+s+s2}{\PYGZdq{}}\PYG{l+s+s2}{heat\PYGZus{}loss\PYGZus{}kw\PYGZus{}per\PYGZus{}c}\PYG{l+s+s2}{\PYGZdq{}}\PYG{p}{,} \PYG{n}{positive}\PYG{o}{=}\PYG{k+kc}{True}\PYG{p}{)}
    \PYG{n}{c} \PYG{o}{=} \PYG{n}{\PYGZus{}number}\PYG{p}{(}\PYG{n}{thermal\PYGZus{}capacity\PYGZus{}kwh\PYGZus{}per\PYGZus{}c}\PYG{p}{,} \PYG{l+s+s2}{\PYGZdq{}}\PYG{l+s+s2}{thermal\PYGZus{}capacity\PYGZus{}kwh\PYGZus{}per\PYGZus{}c}\PYG{l+s+s2}{\PYGZdq{}}\PYG{p}{,} \PYG{n}{positive}\PYG{o}{=}\PYG{k+kc}{True}\PYG{p}{)}
    \PYG{n}{dt} \PYG{o}{=} \PYG{n}{\PYGZus{}number}\PYG{p}{(}\PYG{n}{dt\PYGZus{}hours}\PYG{p}{,} \PYG{l+s+s2}{\PYGZdq{}}\PYG{l+s+s2}{dt\PYGZus{}hours}\PYG{l+s+s2}{\PYGZdq{}}\PYG{p}{,} \PYG{n}{positive}\PYG{o}{=}\PYG{k+kc}{True}\PYG{p}{)}
    \PYG{n}{decay} \PYG{o}{=} \PYG{n}{h} \PYG{o}{*} \PYG{p}{(}\PYG{n}{dt} \PYG{o}{/} \PYG{n}{c}\PYG{p}{)}
    \PYG{n}{a}\PYG{p}{,} \PYG{n}{b} \PYG{o}{=} \PYG{n}{exp}\PYG{p}{(}\PYG{o}{\PYGZhy{}}\PYG{n}{decay}\PYG{p}{)}\PYG{p}{,} \PYG{o}{\PYGZhy{}}\PYG{n}{expm1}\PYG{p}{(}\PYG{o}{\PYGZhy{}}\PYG{n}{decay}\PYG{p}{)} \PYG{o}{/} \PYG{n}{h}
    \PYG{k}{if} \PYG{o+ow}{not} \PYG{n}{isfinite}\PYG{p}{(}\PYG{n}{b}\PYG{p}{)} \PYG{o+ow}{or} \PYG{n}{b} \PYG{o}{\PYGZlt{}}\PYG{o}{=} \PYG{l+m+mf}{1e\PYGZhy{}12}\PYG{p}{:}
        \PYG{k}{raise} \PYG{n+ne}{ValueError}\PYG{p}{(}\PYG{l+s+s2}{\PYGZdq{}}\PYG{l+s+s2}{RC coefficient b is outside LinearMILP numerical resolution; rescale units}\PYG{l+s+s2}{\PYGZdq{}}\PYG{p}{)}
    \PYG{k}{return} \PYG{n}{a}\PYG{p}{,} \PYG{n}{b}
\end{Verbatim}

\endgroup

\Needspace{14\baselineskip}
\subsection*{3．节点热平衡}
按节点、温位和小时组装热源、双向输配、储热与需求项；损失、需求响应和建筑交付量随配置进入等式。变量注册及其他约束见完整文件。

\begingroup
\setstretch{1.0}
\begin{Verbatim}[commandchars=\\\{\},numbers=left,firstnumber=1603,stepnumber=1,fontsize=\fontsize{9.5}{11.5}\selectfont,breaklines=true,breakanywhere=true,frame=single,framesep=3mm,numbersep=5pt,rulecolor=\color{black},xleftmargin=18pt]
\PYG{c+c1}{\PYGZsh{} Node\PYGZhy{}grade\PYGZhy{}time heat balances.}
\PYG{k}{for} \PYG{n}{node\PYGZus{}name}\PYG{p}{,} \PYG{n}{node} \PYG{o+ow}{in} \PYG{n+nb+bp}{self}\PYG{o}{.}\PYG{n}{nodes}\PYG{o}{.}\PYG{n}{items}\PYG{p}{(}\PYG{p}{)}\PYG{p}{:}
    \PYG{k}{for} \PYG{n}{grade\PYGZus{}name}\PYG{p}{,} \PYG{n}{grade} \PYG{o+ow}{in} \PYG{n+nb+bp}{self}\PYG{o}{.}\PYG{n}{grades}\PYG{o}{.}\PYG{n}{items}\PYG{p}{(}\PYG{p}{)}\PYG{p}{:}
        \PYG{n}{demand\PYGZus{}column} \PYG{o}{=} \PYG{n}{node}\PYG{o}{.}\PYG{n}{get}\PYG{p}{(}\PYG{l+s+s2}{\PYGZdq{}}\PYG{l+s+s2}{demand\PYGZus{}columns}\PYG{l+s+s2}{\PYGZdq{}}\PYG{p}{,} \PYG{p}{\PYGZob{}}\PYG{p}{\PYGZcb{}}\PYG{p}{)}\PYG{o}{.}\PYG{n}{get}\PYG{p}{(}\PYG{n}{grade\PYGZus{}name}\PYG{p}{)}
        \PYG{k}{for} \PYG{n}{t} \PYG{o+ow}{in} \PYG{n+nb+bp}{self}\PYG{o}{.}\PYG{n}{times}\PYG{p}{:}
            \PYG{n}{row}\PYG{p}{:} \PYG{n+nb}{dict}\PYG{p}{[}\PYG{n+nb}{int}\PYG{p}{,} \PYG{n+nb}{float}\PYG{p}{]} \PYG{o}{=} \PYG{p}{\PYGZob{}}\PYG{p}{\PYGZcb{}}
            \PYG{k}{for} \PYG{n}{s\PYGZus{}name}\PYG{p}{,} \PYG{n}{source} \PYG{o+ow}{in} \PYG{n+nb+bp}{self}\PYG{o}{.}\PYG{n}{sources}\PYG{o}{.}\PYG{n}{items}\PYG{p}{(}\PYG{p}{)}\PYG{p}{:}
                \PYG{n}{idx} \PYG{o}{=} \PYG{n+nb+bp}{self}\PYG{o}{.}\PYG{n}{variables}\PYG{p}{[}\PYG{l+s+s2}{\PYGZdq{}}\PYG{l+s+s2}{source\PYGZus{}heat}\PYG{l+s+s2}{\PYGZdq{}}\PYG{p}{]}\PYG{o}{.}\PYG{n}{get}\PYG{p}{(}\PYG{p}{(}\PYG{n}{s\PYGZus{}name}\PYG{p}{,} \PYG{n}{grade\PYGZus{}name}\PYG{p}{,} \PYG{n}{t}\PYG{p}{)}\PYG{p}{)}
                \PYG{k}{if} \PYG{n}{idx} \PYG{o+ow}{is} \PYG{o+ow}{not} \PYG{k+kc}{None} \PYG{o+ow}{and} \PYG{n}{source}\PYG{p}{[}\PYG{l+s+s2}{\PYGZdq{}}\PYG{l+s+s2}{node}\PYG{l+s+s2}{\PYGZdq{}}\PYG{p}{]} \PYG{o}{==} \PYG{n}{node\PYGZus{}name}\PYG{p}{:}
                    \PYG{n+nb+bp}{self}\PYG{o}{.}\PYG{n}{\PYGZus{}add}\PYG{p}{(}\PYG{n}{row}\PYG{p}{,} \PYG{n}{idx}\PYG{p}{,} \PYG{l+m+mf}{1.0}\PYG{p}{)}
            \PYG{k}{for} \PYG{n}{e\PYGZus{}name}\PYG{p}{,} \PYG{n}{pipe} \PYG{o+ow}{in} \PYG{n+nb+bp}{self}\PYG{o}{.}\PYG{n}{pipes}\PYG{o}{.}\PYG{n}{items}\PYG{p}{(}\PYG{p}{)}\PYG{p}{:}
                \PYG{n}{forward} \PYG{o}{=} \PYG{n+nb+bp}{self}\PYG{o}{.}\PYG{n}{variables}\PYG{p}{[}\PYG{l+s+s2}{\PYGZdq{}}\PYG{l+s+s2}{pipe\PYGZus{}flow\PYGZus{}forward}\PYG{l+s+s2}{\PYGZdq{}}\PYG{p}{]}\PYG{p}{[}\PYG{p}{(}\PYG{n}{e\PYGZus{}name}\PYG{p}{,} \PYG{n}{grade\PYGZus{}name}\PYG{p}{,} \PYG{n}{t}\PYG{p}{)}\PYG{p}{]}
                \PYG{k}{if} \PYG{n}{pipe}\PYG{p}{[}\PYG{l+s+s2}{\PYGZdq{}}\PYG{l+s+s2}{from}\PYG{l+s+s2}{\PYGZdq{}}\PYG{p}{]} \PYG{o}{==} \PYG{n}{node\PYGZus{}name}\PYG{p}{:}
                    \PYG{n+nb+bp}{self}\PYG{o}{.}\PYG{n}{\PYGZus{}add}\PYG{p}{(}\PYG{n}{row}\PYG{p}{,} \PYG{n}{forward}\PYG{p}{,} \PYG{o}{\PYGZhy{}}\PYG{l+m+mf}{1.0}\PYG{p}{)}
                \PYG{k}{if} \PYG{n}{pipe}\PYG{p}{[}\PYG{l+s+s2}{\PYGZdq{}}\PYG{l+s+s2}{to}\PYG{l+s+s2}{\PYGZdq{}}\PYG{p}{]} \PYG{o}{==} \PYG{n}{node\PYGZus{}name}\PYG{p}{:}
                    \PYG{n+nb+bp}{self}\PYG{o}{.}\PYG{n}{\PYGZus{}add}\PYG{p}{(}\PYG{n}{row}\PYG{p}{,} \PYG{n}{forward}\PYG{p}{,} \PYG{n+nb+bp}{self}\PYG{o}{.}\PYG{n}{\PYGZus{}pipe\PYGZus{}efficiency}\PYG{p}{(}\PYG{n}{pipe}\PYG{p}{,} \PYG{n}{t}\PYG{p}{)}\PYG{p}{)}
                    \PYG{n}{activated} \PYG{o}{=} \PYG{n+nb+bp}{self}\PYG{o}{.}\PYG{n}{variables}\PYG{p}{[}\PYG{l+s+s2}{\PYGZdq{}}\PYG{l+s+s2}{pipe\PYGZus{}retrofit\PYGZus{}flow\PYGZus{}forward}\PYG{l+s+s2}{\PYGZdq{}}\PYG{p}{]}\PYG{o}{.}\PYG{n}{get}\PYG{p}{(}
                        \PYG{p}{(}\PYG{n}{e\PYGZus{}name}\PYG{p}{,} \PYG{n}{grade\PYGZus{}name}\PYG{p}{,} \PYG{n}{t}\PYG{p}{)}
                    \PYG{p}{)}
                    \PYG{k}{if} \PYG{n}{activated} \PYG{o+ow}{is} \PYG{o+ow}{not} \PYG{k+kc}{None}\PYG{p}{:}
                        \PYG{n}{gain} \PYG{o}{=} \PYG{n+nb+bp}{self}\PYG{o}{.}\PYG{n}{\PYGZus{}pipe\PYGZus{}retrofit\PYGZus{}efficiency}\PYG{p}{(}\PYG{n}{pipe}\PYG{p}{,} \PYG{n}{t}\PYG{p}{)} \PYG{o}{\PYGZhy{}} \PYG{n+nb+bp}{self}\PYG{o}{.}\PYG{n}{\PYGZus{}pipe\PYGZus{}efficiency}\PYG{p}{(}\PYG{n}{pipe}\PYG{p}{,} \PYG{n}{t}\PYG{p}{)}
                        \PYG{n+nb+bp}{self}\PYG{o}{.}\PYG{n}{\PYGZus{}add}\PYG{p}{(}\PYG{n}{row}\PYG{p}{,} \PYG{n}{activated}\PYG{p}{,} \PYG{n+nb}{max}\PYG{p}{(}\PYG{l+m+mf}{0.0}\PYG{p}{,} \PYG{n}{gain}\PYG{p}{)}\PYG{p}{)}
                \PYG{n}{reverse} \PYG{o}{=} \PYG{n+nb+bp}{self}\PYG{o}{.}\PYG{n}{variables}\PYG{p}{[}\PYG{l+s+s2}{\PYGZdq{}}\PYG{l+s+s2}{pipe\PYGZus{}flow\PYGZus{}reverse}\PYG{l+s+s2}{\PYGZdq{}}\PYG{p}{]}\PYG{o}{.}\PYG{n}{get}\PYG{p}{(}\PYG{p}{(}\PYG{n}{e\PYGZus{}name}\PYG{p}{,} \PYG{n}{grade\PYGZus{}name}\PYG{p}{,} \PYG{n}{t}\PYG{p}{)}\PYG{p}{)}
                \PYG{k}{if} \PYG{n}{reverse} \PYG{o+ow}{is} \PYG{o+ow}{not} \PYG{k+kc}{None}\PYG{p}{:}
                    \PYG{k}{if} \PYG{n}{pipe}\PYG{p}{[}\PYG{l+s+s2}{\PYGZdq{}}\PYG{l+s+s2}{to}\PYG{l+s+s2}{\PYGZdq{}}\PYG{p}{]} \PYG{o}{==} \PYG{n}{node\PYGZus{}name}\PYG{p}{:}
                        \PYG{n+nb+bp}{self}\PYG{o}{.}\PYG{n}{\PYGZus{}add}\PYG{p}{(}\PYG{n}{row}\PYG{p}{,} \PYG{n}{reverse}\PYG{p}{,} \PYG{o}{\PYGZhy{}}\PYG{l+m+mf}{1.0}\PYG{p}{)}
                    \PYG{k}{if} \PYG{n}{pipe}\PYG{p}{[}\PYG{l+s+s2}{\PYGZdq{}}\PYG{l+s+s2}{from}\PYG{l+s+s2}{\PYGZdq{}}\PYG{p}{]} \PYG{o}{==} \PYG{n}{node\PYGZus{}name}\PYG{p}{:}
                        \PYG{n+nb+bp}{self}\PYG{o}{.}\PYG{n}{\PYGZus{}add}\PYG{p}{(}\PYG{n}{row}\PYG{p}{,} \PYG{n}{reverse}\PYG{p}{,} \PYG{n+nb+bp}{self}\PYG{o}{.}\PYG{n}{\PYGZus{}pipe\PYGZus{}efficiency}\PYG{p}{(}\PYG{n}{pipe}\PYG{p}{,} \PYG{n}{t}\PYG{p}{)}\PYG{p}{)}
                        \PYG{n}{activated\PYGZus{}reverse} \PYG{o}{=} \PYG{n+nb+bp}{self}\PYG{o}{.}\PYG{n}{variables}\PYG{p}{[}\PYG{l+s+s2}{\PYGZdq{}}\PYG{l+s+s2}{pipe\PYGZus{}retrofit\PYGZus{}flow\PYGZus{}reverse}\PYG{l+s+s2}{\PYGZdq{}}\PYG{p}{]}\PYG{o}{.}\PYG{n}{get}\PYG{p}{(}
                            \PYG{p}{(}\PYG{n}{e\PYGZus{}name}\PYG{p}{,} \PYG{n}{grade\PYGZus{}name}\PYG{p}{,} \PYG{n}{t}\PYG{p}{)}
                        \PYG{p}{)}
                        \PYG{k}{if} \PYG{n}{activated\PYGZus{}reverse} \PYG{o+ow}{is} \PYG{o+ow}{not} \PYG{k+kc}{None}\PYG{p}{:}
                            \PYG{n}{gain} \PYG{o}{=} \PYG{n+nb+bp}{self}\PYG{o}{.}\PYG{n}{\PYGZus{}pipe\PYGZus{}retrofit\PYGZus{}efficiency}\PYG{p}{(}\PYG{n}{pipe}\PYG{p}{,} \PYG{n}{t}\PYG{p}{)} \PYG{o}{\PYGZhy{}} \PYG{n+nb+bp}{self}\PYG{o}{.}\PYG{n}{\PYGZus{}pipe\PYGZus{}efficiency}\PYG{p}{(}\PYG{n}{pipe}\PYG{p}{,} \PYG{n}{t}\PYG{p}{)}
                            \PYG{n+nb+bp}{self}\PYG{o}{.}\PYG{n}{\PYGZus{}add}\PYG{p}{(}\PYG{n}{row}\PYG{p}{,} \PYG{n}{activated\PYGZus{}reverse}\PYG{p}{,} \PYG{n+nb}{max}\PYG{p}{(}\PYG{l+m+mf}{0.0}\PYG{p}{,} \PYG{n}{gain}\PYG{p}{)}\PYG{p}{)}
            \PYG{k}{for} \PYG{n}{st\PYGZus{}name}\PYG{p}{,} \PYG{n}{storage} \PYG{o+ow}{in} \PYG{n+nb+bp}{self}\PYG{o}{.}\PYG{n}{storages}\PYG{o}{.}\PYG{n}{items}\PYG{p}{(}\PYG{p}{)}\PYG{p}{:}
                \PYG{k}{if} \PYG{n}{storage}\PYG{p}{[}\PYG{l+s+s2}{\PYGZdq{}}\PYG{l+s+s2}{node}\PYG{l+s+s2}{\PYGZdq{}}\PYG{p}{]} \PYG{o}{==} \PYG{n}{node\PYGZus{}name} \PYG{o+ow}{and} \PYG{n}{storage}\PYG{p}{[}\PYG{l+s+s2}{\PYGZdq{}}\PYG{l+s+s2}{grade}\PYG{l+s+s2}{\PYGZdq{}}\PYG{p}{]} \PYG{o}{==} \PYG{n}{grade\PYGZus{}name}\PYG{p}{:}
                    \PYG{n+nb+bp}{self}\PYG{o}{.}\PYG{n}{\PYGZus{}add}\PYG{p}{(}\PYG{n}{row}\PYG{p}{,} \PYG{n+nb+bp}{self}\PYG{o}{.}\PYG{n}{variables}\PYG{p}{[}\PYG{l+s+s2}{\PYGZdq{}}\PYG{l+s+s2}{storage\PYGZus{}discharge}\PYG{l+s+s2}{\PYGZdq{}}\PYG{p}{]}\PYG{p}{[}\PYG{p}{(}\PYG{n}{st\PYGZus{}name}\PYG{p}{,} \PYG{n}{t}\PYG{p}{)}\PYG{p}{]}\PYG{p}{,} \PYG{l+m+mf}{1.0}\PYG{p}{)}
                    \PYG{n+nb+bp}{self}\PYG{o}{.}\PYG{n}{\PYGZus{}add}\PYG{p}{(}\PYG{n}{row}\PYG{p}{,} \PYG{n+nb+bp}{self}\PYG{o}{.}\PYG{n}{variables}\PYG{p}{[}\PYG{l+s+s2}{\PYGZdq{}}\PYG{l+s+s2}{storage\PYGZus{}charge}\PYG{l+s+s2}{\PYGZdq{}}\PYG{p}{]}\PYG{p}{[}\PYG{p}{(}\PYG{n}{st\PYGZus{}name}\PYG{p}{,} \PYG{n}{t}\PYG{p}{)}\PYG{p}{]}\PYG{p}{,} \PYG{o}{\PYGZhy{}}\PYG{l+m+mf}{1.0}\PYG{p}{)}
            \PYG{n}{unmet} \PYG{o}{=} \PYG{n+nb+bp}{self}\PYG{o}{.}\PYG{n}{variables}\PYG{p}{[}\PYG{l+s+s2}{\PYGZdq{}}\PYG{l+s+s2}{unserved}\PYG{l+s+s2}{\PYGZdq{}}\PYG{p}{]}\PYG{o}{.}\PYG{n}{get}\PYG{p}{(}\PYG{p}{(}\PYG{n}{node\PYGZus{}name}\PYG{p}{,} \PYG{n}{grade\PYGZus{}name}\PYG{p}{,} \PYG{n}{t}\PYG{p}{)}\PYG{p}{)}
            \PYG{k}{if} \PYG{n}{unmet} \PYG{o+ow}{is} \PYG{o+ow}{not} \PYG{k+kc}{None}\PYG{p}{:}
                \PYG{n+nb+bp}{self}\PYG{o}{.}\PYG{n}{\PYGZus{}add}\PYG{p}{(}\PYG{n}{row}\PYG{p}{,} \PYG{n}{unmet}\PYG{p}{,} \PYG{l+m+mf}{1.0}\PYG{p}{)}
            \PYG{n}{up} \PYG{o}{=} \PYG{n+nb+bp}{self}\PYG{o}{.}\PYG{n}{variables}\PYG{p}{[}\PYG{l+s+s2}{\PYGZdq{}}\PYG{l+s+s2}{demand\PYGZus{}up}\PYG{l+s+s2}{\PYGZdq{}}\PYG{p}{]}\PYG{o}{.}\PYG{n}{get}\PYG{p}{(}\PYG{p}{(}\PYG{n}{node\PYGZus{}name}\PYG{p}{,} \PYG{n}{grade\PYGZus{}name}\PYG{p}{,} \PYG{n}{t}\PYG{p}{)}\PYG{p}{)}
            \PYG{n}{down} \PYG{o}{=} \PYG{n+nb+bp}{self}\PYG{o}{.}\PYG{n}{variables}\PYG{p}{[}\PYG{l+s+s2}{\PYGZdq{}}\PYG{l+s+s2}{demand\PYGZus{}down}\PYG{l+s+s2}{\PYGZdq{}}\PYG{p}{]}\PYG{o}{.}\PYG{n}{get}\PYG{p}{(}\PYG{p}{(}\PYG{n}{node\PYGZus{}name}\PYG{p}{,} \PYG{n}{grade\PYGZus{}name}\PYG{p}{,} \PYG{n}{t}\PYG{p}{)}\PYG{p}{)}
            \PYG{k}{if} \PYG{n}{up} \PYG{o+ow}{is} \PYG{o+ow}{not} \PYG{k+kc}{None}\PYG{p}{:}
                \PYG{n+nb+bp}{self}\PYG{o}{.}\PYG{n}{\PYGZus{}add}\PYG{p}{(}\PYG{n}{row}\PYG{p}{,} \PYG{n}{up}\PYG{p}{,} \PYG{o}{\PYGZhy{}}\PYG{l+m+mf}{1.0}\PYG{p}{)}
            \PYG{k}{if} \PYG{n}{down} \PYG{o+ow}{is} \PYG{o+ow}{not} \PYG{k+kc}{None}\PYG{p}{:}
                \PYG{n+nb+bp}{self}\PYG{o}{.}\PYG{n}{\PYGZus{}add}\PYG{p}{(}\PYG{n}{row}\PYG{p}{,} \PYG{n}{down}\PYG{p}{,} \PYG{l+m+mf}{1.0}\PYG{p}{)}
            \PYG{k}{if} \PYG{n+nb+bp}{self}\PYG{o}{.}\PYG{n}{building\PYGZus{}thermal} \PYG{o+ow}{is} \PYG{o+ow}{not} \PYG{k+kc}{None} \PYG{o+ow}{and} \PYG{p}{(}\PYG{n}{node\PYGZus{}name}\PYG{p}{,}\PYG{n}{grade\PYGZus{}name}\PYG{p}{)} \PYG{o+ow}{in} \PYG{n+nb+bp}{self}\PYG{o}{.}\PYG{n}{building\PYGZus{}thermal}\PYG{o}{.}\PYG{n}{covered\PYGZus{}node\PYGZus{}grades}\PYG{p}{:}
                \PYG{n}{delivered} \PYG{o}{=} \PYG{n+nb+bp}{self}\PYG{o}{.}\PYG{n}{building\PYGZus{}thermal}\PYG{o}{.}\PYG{n}{node\PYGZus{}heat}\PYG{p}{[}\PYG{n}{node\PYGZus{}name}\PYG{p}{,}\PYG{n}{grade\PYGZus{}name}\PYG{p}{,}\PYG{n}{t}\PYG{p}{]}
                \PYG{n+nb+bp}{self}\PYG{o}{.}\PYG{n}{\PYGZus{}add}\PYG{p}{(}\PYG{n}{row}\PYG{p}{,}\PYG{n}{delivered}\PYG{p}{,}\PYG{o}{\PYGZhy{}}\PYG{l+m+mf}{1.0}\PYG{p}{)}
                \PYG{n}{p}\PYG{o}{.}\PYG{n}{add\PYGZus{}constraint}\PYG{p}{(}\PYG{n}{row}\PYG{p}{,}\PYG{n}{lb}\PYG{o}{=}\PYG{l+m+mf}{0.0}\PYG{p}{,}\PYG{n}{ub}\PYG{o}{=}\PYG{l+m+mf}{0.0}\PYG{p}{)}
                \PYG{n+nb+bp}{self}\PYG{o}{.}\PYG{n}{\PYGZus{}record\PYGZus{}time\PYGZus{}metric}\PYG{p}{(}\PYG{l+s+s2}{\PYGZdq{}}\PYG{l+s+s2}{demand\PYGZus{}kwh}\PYG{l+s+s2}{\PYGZdq{}}\PYG{p}{,}\PYG{n}{delivered}\PYG{p}{,}\PYG{l+m+mf}{1.0}\PYG{p}{,}\PYG{n}{t}\PYG{p}{)}
                \PYG{n+nb+bp}{self}\PYG{o}{.}\PYG{n}{\PYGZus{}record\PYGZus{}time\PYGZus{}metric}\PYG{p}{(}\PYG{l+s+s2}{\PYGZdq{}}\PYG{l+s+s2}{building\PYGZus{}heat\PYGZus{}delivered\PYGZus{}kwh}\PYG{l+s+s2}{\PYGZdq{}}\PYG{p}{,}\PYG{n}{delivered}\PYG{p}{,}\PYG{l+m+mf}{1.0}\PYG{p}{,}\PYG{n}{t}\PYG{p}{)}
                \PYG{n+nb+bp}{self}\PYG{o}{.}\PYG{n}{\PYGZus{}record\PYGZus{}time\PYGZus{}metric}\PYG{p}{(}\PYG{l+s+s2}{\PYGZdq{}}\PYG{l+s+s2}{demand\PYGZus{}exergy\PYGZus{}kwh}\PYG{l+s+s2}{\PYGZdq{}}\PYG{p}{,}\PYG{n}{delivered}\PYG{p}{,}\PYG{n+nb+bp}{self}\PYG{o}{.}\PYG{n}{\PYGZus{}exergy\PYGZus{}factor}\PYG{p}{(}\PYG{n+nb}{float}\PYG{p}{(}\PYG{n}{grade}\PYG{p}{[}\PYG{l+s+s2}{\PYGZdq{}}\PYG{l+s+s2}{temperature\PYGZus{}c}\PYG{l+s+s2}{\PYGZdq{}}\PYG{p}{]}\PYG{p}{)}\PYG{p}{)}\PYG{p}{,}\PYG{n}{t}\PYG{p}{)}
                \PYG{n}{reference} \PYG{o}{=} \PYG{n+nb+bp}{self}\PYG{o}{.}\PYG{n}{\PYGZus{}demand\PYGZus{}value}\PYG{p}{(}\PYG{n}{node}\PYG{p}{,}\PYG{n}{demand\PYGZus{}column}\PYG{p}{,}\PYG{n}{t}\PYG{p}{)} \PYG{k}{if} \PYG{n}{demand\PYGZus{}column} \PYG{k}{else} \PYG{l+m+mf}{0.0}
                \PYG{n+nb+bp}{self}\PYG{o}{.}\PYG{n}{metric\PYGZus{}constants}\PYG{p}{[}\PYG{l+s+s2}{\PYGZdq{}}\PYG{l+s+s2}{building\PYGZus{}reference\PYGZus{}heat\PYGZus{}kwh}\PYG{l+s+s2}{\PYGZdq{}}\PYG{p}{]} \PYG{o}{+}\PYG{o}{=} \PYG{n}{reference}\PYG{o}{*}\PYG{n+nb+bp}{self}\PYG{o}{.}\PYG{n}{\PYGZus{}op\PYGZus{}weight}\PYG{p}{(}\PYG{n}{t}\PYG{p}{)}
                \PYG{n+nb+bp}{self}\PYG{o}{.}\PYG{n}{scenario\PYGZus{}metric\PYGZus{}constants}\PYG{p}{[}\PYG{n+nb+bp}{self}\PYG{o}{.}\PYG{n}{\PYGZus{}scenario}\PYG{p}{(}\PYG{n}{t}\PYG{p}{)}\PYG{p}{]}\PYG{p}{[}\PYG{l+s+s2}{\PYGZdq{}}\PYG{l+s+s2}{building\PYGZus{}reference\PYGZus{}heat\PYGZus{}kwh}\PYG{l+s+s2}{\PYGZdq{}}\PYG{p}{]} \PYG{o}{+}\PYG{o}{=} \PYG{n}{reference}\PYG{o}{*}\PYG{n+nb+bp}{self}\PYG{o}{.}\PYG{n}{\PYGZus{}physical\PYGZus{}weight}\PYG{p}{(}\PYG{n}{t}\PYG{p}{)}
                \PYG{n+nb+bp}{self}\PYG{o}{.}\PYG{n}{\PYGZus{}record\PYGZus{}metering\PYGZus{}reference}\PYG{p}{(}\PYG{n}{node\PYGZus{}name}\PYG{p}{,} \PYG{n}{grade}\PYG{p}{,} \PYG{n}{reference}\PYG{p}{,} \PYG{n}{t}\PYG{p}{)}
                \PYG{k}{continue}
            \PYG{n}{demand} \PYG{o}{=} \PYG{n+nb+bp}{self}\PYG{o}{.}\PYG{n}{\PYGZus{}demand\PYGZus{}value}\PYG{p}{(}\PYG{n}{node}\PYG{p}{,} \PYG{n}{demand\PYGZus{}column}\PYG{p}{,} \PYG{n}{t}\PYG{p}{)} \PYG{k}{if} \PYG{n}{demand\PYGZus{}column} \PYG{k}{else} \PYG{l+m+mf}{0.0}
            \PYG{n}{metering} \PYG{o}{=} \PYG{n}{node}\PYG{o}{.}\PYG{n}{get}\PYG{p}{(}\PYG{l+s+s2}{\PYGZdq{}}\PYG{l+s+s2}{heat\PYGZus{}metering}\PYG{l+s+s2}{\PYGZdq{}}\PYG{p}{,} \PYG{p}{\PYGZob{}}\PYG{p}{\PYGZcb{}}\PYG{p}{)}
            \PYG{n}{existing\PYGZus{}metering} \PYG{o}{=} \PYG{n+nb}{float}\PYG{p}{(}
                \PYG{n}{metering}\PYG{o}{.}\PYG{n}{get}\PYG{p}{(}\PYG{l+s+s2}{\PYGZdq{}}\PYG{l+s+s2}{existing\PYGZus{}coverage\PYGZus{}fraction}\PYG{l+s+s2}{\PYGZdq{}}\PYG{p}{,} \PYG{l+m+mf}{0.0}\PYG{p}{)}
            \PYG{p}{)}
            \PYG{n}{saving\PYGZus{}fraction} \PYG{o}{=} \PYG{n+nb}{float}\PYG{p}{(}
                \PYG{n}{metering}\PYG{o}{.}\PYG{n}{get}\PYG{p}{(}\PYG{l+s+s2}{\PYGZdq{}}\PYG{l+s+s2}{saving\PYGZus{}fraction\PYGZus{}at\PYGZus{}full\PYGZus{}coverage}\PYG{l+s+s2}{\PYGZdq{}}\PYG{p}{,} \PYG{l+m+mf}{0.0}\PYG{p}{)}
            \PYG{p}{)}
            \PYG{n}{existing\PYGZus{}metering\PYGZus{}saving} \PYG{o}{=} \PYG{p}{(}
                \PYG{n}{demand} \PYG{o}{*} \PYG{n}{existing\PYGZus{}metering} \PYG{o}{*} \PYG{n}{saving\PYGZus{}fraction}
            \PYG{p}{)}
            \PYG{n}{metering\PYGZus{}idx} \PYG{o}{=} \PYG{n+nb+bp}{self}\PYG{o}{.}\PYG{n}{variables}\PYG{p}{[}\PYG{l+s+s2}{\PYGZdq{}}\PYG{l+s+s2}{heat\PYGZus{}metering\PYGZus{}coverage}\PYG{l+s+s2}{\PYGZdq{}}\PYG{p}{]}\PYG{o}{.}\PYG{n}{get}\PYG{p}{(}
                \PYG{n}{node\PYGZus{}name}
            \PYG{p}{)}
            \PYG{k}{if} \PYG{n}{metering\PYGZus{}idx} \PYG{o+ow}{is} \PYG{o+ow}{not} \PYG{k+kc}{None} \PYG{o+ow}{and} \PYG{n}{demand} \PYG{o}{\PYGZgt{}} \PYG{l+m+mf}{0.0} \PYG{o+ow}{and} \PYG{n}{saving\PYGZus{}fraction} \PYG{o}{\PYGZgt{}} \PYG{l+m+mf}{0.0}\PYG{p}{:}
                \PYG{n+nb+bp}{self}\PYG{o}{.}\PYG{n}{\PYGZus{}add}\PYG{p}{(}\PYG{n}{row}\PYG{p}{,} \PYG{n}{metering\PYGZus{}idx}\PYG{p}{,} \PYG{n}{demand} \PYG{o}{*} \PYG{n}{saving\PYGZus{}fraction}\PYG{p}{)}
            \PYG{n}{retrofit} \PYG{o}{=} \PYG{n}{node}\PYG{o}{.}\PYG{n}{get}\PYG{p}{(}\PYG{l+s+s2}{\PYGZdq{}}\PYG{l+s+s2}{temperature\PYGZus{}retrofit}\PYG{l+s+s2}{\PYGZdq{}}\PYG{p}{)}
            \PYG{n}{retrofit\PYGZus{}idx} \PYG{o}{=} \PYG{n+nb+bp}{self}\PYG{o}{.}\PYG{n}{variables}\PYG{p}{[}\PYG{l+s+s2}{\PYGZdq{}}\PYG{l+s+s2}{terminal\PYGZus{}temperature\PYGZus{}retrofit}\PYG{l+s+s2}{\PYGZdq{}}\PYG{p}{]}\PYG{o}{.}\PYG{n}{get}\PYG{p}{(}\PYG{n}{node\PYGZus{}name}\PYG{p}{)}
            \PYG{k}{if} \PYG{n}{retrofit} \PYG{o+ow}{and} \PYG{n}{retrofit\PYGZus{}idx} \PYG{o+ow}{is} \PYG{o+ow}{not} \PYG{k+kc}{None}\PYG{p}{:}
                \PYG{n}{source\PYGZus{}grade} \PYG{o}{=} \PYG{n+nb}{str}\PYG{p}{(}\PYG{n}{retrofit}\PYG{p}{[}\PYG{l+s+s2}{\PYGZdq{}}\PYG{l+s+s2}{source\PYGZus{}grade}\PYG{l+s+s2}{\PYGZdq{}}\PYG{p}{]}\PYG{p}{)}
                \PYG{n}{target\PYGZus{}grade} \PYG{o}{=} \PYG{n+nb}{str}\PYG{p}{(}\PYG{n}{retrofit}\PYG{p}{[}\PYG{l+s+s2}{\PYGZdq{}}\PYG{l+s+s2}{target\PYGZus{}grade}\PYG{l+s+s2}{\PYGZdq{}}\PYG{p}{]}\PYG{p}{)}
                \PYG{n}{source\PYGZus{}column} \PYG{o}{=} \PYG{n}{node}\PYG{p}{[}\PYG{l+s+s2}{\PYGZdq{}}\PYG{l+s+s2}{demand\PYGZus{}columns}\PYG{l+s+s2}{\PYGZdq{}}\PYG{p}{]}\PYG{p}{[}\PYG{n}{source\PYGZus{}grade}\PYG{p}{]}
                \PYG{n}{shift} \PYG{o}{=} \PYG{p}{(}
                    \PYG{n+nb}{float}\PYG{p}{(}\PYG{n}{retrofit}\PYG{o}{.}\PYG{n}{get}\PYG{p}{(}\PYG{l+s+s2}{\PYGZdq{}}\PYG{l+s+s2}{eligible\PYGZus{}fraction}\PYG{l+s+s2}{\PYGZdq{}}\PYG{p}{,} \PYG{l+m+mf}{1.0}\PYG{p}{)}\PYG{p}{)}
                    \PYG{o}{*} \PYG{n+nb+bp}{self}\PYG{o}{.}\PYG{n}{\PYGZus{}demand\PYGZus{}value}\PYG{p}{(}\PYG{n}{node}\PYG{p}{,} \PYG{n}{source\PYGZus{}column}\PYG{p}{,} \PYG{n}{t}\PYG{p}{)}
                \PYG{p}{)}
                \PYG{k}{if} \PYG{n}{grade\PYGZus{}name} \PYG{o}{==} \PYG{n}{source\PYGZus{}grade}\PYG{p}{:}
                    \PYG{n+nb+bp}{self}\PYG{o}{.}\PYG{n}{\PYGZus{}add}\PYG{p}{(}\PYG{n}{row}\PYG{p}{,} \PYG{n}{retrofit\PYGZus{}idx}\PYG{p}{,} \PYG{n}{shift}\PYG{p}{)}
                \PYG{k}{elif} \PYG{n}{grade\PYGZus{}name} \PYG{o}{==} \PYG{n}{target\PYGZus{}grade}\PYG{p}{:}
                    \PYG{n+nb+bp}{self}\PYG{o}{.}\PYG{n}{\PYGZus{}add}\PYG{p}{(}\PYG{n}{row}\PYG{p}{,} \PYG{n}{retrofit\PYGZus{}idx}\PYG{p}{,} \PYG{o}{\PYGZhy{}}\PYG{n}{shift}\PYG{p}{)}
            \PYG{n}{effective\PYGZus{}constant\PYGZus{}demand} \PYG{o}{=} \PYG{n}{demand} \PYG{o}{\PYGZhy{}} \PYG{n}{existing\PYGZus{}metering\PYGZus{}saving}
            \PYG{n}{p}\PYG{o}{.}\PYG{n}{add\PYGZus{}constraint}\PYG{p}{(}
                \PYG{n}{row}\PYG{p}{,}
                \PYG{n}{lb}\PYG{o}{=}\PYG{n}{effective\PYGZus{}constant\PYGZus{}demand}\PYG{p}{,}
                \PYG{n}{ub}\PYG{o}{=}\PYG{n}{effective\PYGZus{}constant\PYGZus{}demand}\PYG{p}{,}
            \PYG{p}{)}
\end{Verbatim}

\endgroup

\Needspace{14\baselineskip}
\subsection*{4．储热递推与边界}
约束储热容量、充放功率、状态递推与期末边界；严格充放互斥按配置启用。此前的变量注册不在本段。

\begingroup
\setstretch{1.0}
\begin{Verbatim}[commandchars=\\\{\},numbers=left,firstnumber=1425,stepnumber=1,fontsize=\fontsize{9.5}{11.5}\selectfont,breaklines=true,breakanywhere=true,frame=single,framesep=3mm,numbersep=5pt,rulecolor=\color{black},xleftmargin=18pt]
    \PYG{n}{soc\PYGZus{}row} \PYG{o}{=} \PYG{p}{\PYGZob{}}\PYG{n}{soc}\PYG{p}{:} \PYG{l+m+mf}{1.0}\PYG{p}{\PYGZcb{}}
    \PYG{n}{ch\PYGZus{}row} \PYG{o}{=} \PYG{p}{\PYGZob{}}\PYG{n}{ch}\PYG{p}{:} \PYG{l+m+mf}{1.0}\PYG{p}{\PYGZcb{}}
    \PYG{n}{dis\PYGZus{}row} \PYG{o}{=} \PYG{p}{\PYGZob{}}\PYG{n}{dis}\PYG{p}{:} \PYG{l+m+mf}{1.0}\PYG{p}{\PYGZcb{}}
    \PYG{k}{if} \PYG{n}{cap\PYGZus{}idx} \PYG{o+ow}{is} \PYG{o+ow}{not} \PYG{k+kc}{None}\PYG{p}{:}
        \PYG{n}{soc\PYGZus{}row}\PYG{p}{[}\PYG{n}{cap\PYGZus{}idx}\PYG{p}{]} \PYG{o}{=} \PYG{o}{\PYGZhy{}}\PYG{l+m+mf}{1.0}
        \PYG{n}{ch\PYGZus{}row}\PYG{p}{[}\PYG{n}{cap\PYGZus{}idx}\PYG{p}{]} \PYG{o}{=} \PYG{o}{\PYGZhy{}}\PYG{n}{power\PYGZus{}ratio}
        \PYG{n}{dis\PYGZus{}row}\PYG{p}{[}\PYG{n}{cap\PYGZus{}idx}\PYG{p}{]} \PYG{o}{=} \PYG{o}{\PYGZhy{}}\PYG{n}{power\PYGZus{}ratio}
    \PYG{n}{p}\PYG{o}{.}\PYG{n}{add\PYGZus{}constraint}\PYG{p}{(}\PYG{n}{soc\PYGZus{}row}\PYG{p}{,} \PYG{n}{ub}\PYG{o}{=}\PYG{n}{existing}\PYG{p}{)}
    \PYG{n}{p}\PYG{o}{.}\PYG{n}{add\PYGZus{}constraint}\PYG{p}{(}\PYG{n}{ch\PYGZus{}row}\PYG{p}{,} \PYG{n}{ub}\PYG{o}{=}\PYG{n}{power\PYGZus{}ratio} \PYG{o}{*} \PYG{n}{existing}\PYG{p}{)}
    \PYG{n}{p}\PYG{o}{.}\PYG{n}{add\PYGZus{}constraint}\PYG{p}{(}\PYG{n}{dis\PYGZus{}row}\PYG{p}{,} \PYG{n}{ub}\PYG{o}{=}\PYG{n}{power\PYGZus{}ratio} \PYG{o}{*} \PYG{n}{existing}\PYG{p}{)}
    \PYG{k}{if} \PYG{n}{strict\PYGZus{}storage\PYGZus{}mode}\PYG{p}{:}
        \PYG{n}{mode} \PYG{o}{=} \PYG{n}{p}\PYG{o}{.}\PYG{n}{add\PYGZus{}var}\PYG{p}{(}\PYG{l+s+sa}{f}\PYG{l+s+s2}{\PYGZdq{}}\PYG{l+s+s2}{storage\PYGZus{}mode[}\PYG{l+s+si}{\PYGZob{}}\PYG{n}{st\PYGZus{}name}\PYG{l+s+si}{\PYGZcb{}}\PYG{l+s+s2}{,}\PYG{l+s+si}{\PYGZob{}}\PYG{n}{t}\PYG{l+s+si}{\PYGZcb{}}\PYG{l+s+s2}{]}\PYG{l+s+s2}{\PYGZdq{}}\PYG{p}{,} \PYG{n}{ub}\PYG{o}{=}\PYG{l+m+mf}{1.0}\PYG{p}{,} \PYG{n}{integer}\PYG{o}{=}\PYG{k+kc}{True}\PYG{p}{)}
        \PYG{n+nb+bp}{self}\PYG{o}{.}\PYG{n}{variables}\PYG{p}{[}\PYG{l+s+s2}{\PYGZdq{}}\PYG{l+s+s2}{storage\PYGZus{}mode}\PYG{l+s+s2}{\PYGZdq{}}\PYG{p}{]}\PYG{p}{[}\PYG{p}{(}\PYG{n}{st\PYGZus{}name}\PYG{p}{,} \PYG{n}{t}\PYG{p}{)}\PYG{p}{]} \PYG{o}{=} \PYG{n}{mode}
        \PYG{n}{p}\PYG{o}{.}\PYG{n}{add\PYGZus{}constraint}\PYG{p}{(}\PYG{p}{\PYGZob{}}\PYG{n}{ch}\PYG{p}{:} \PYG{l+m+mf}{1.0}\PYG{p}{,} \PYG{n}{mode}\PYG{p}{:} \PYG{o}{\PYGZhy{}}\PYG{n}{max\PYGZus{}power}\PYG{p}{\PYGZcb{}}\PYG{p}{,} \PYG{n}{ub}\PYG{o}{=}\PYG{l+m+mf}{0.0}\PYG{p}{)}
        \PYG{n}{p}\PYG{o}{.}\PYG{n}{add\PYGZus{}constraint}\PYG{p}{(}\PYG{p}{\PYGZob{}}\PYG{n}{dis}\PYG{p}{:} \PYG{l+m+mf}{1.0}\PYG{p}{,} \PYG{n}{mode}\PYG{p}{:} \PYG{n}{max\PYGZus{}power}\PYG{p}{\PYGZcb{}}\PYG{p}{,} \PYG{n}{ub}\PYG{o}{=}\PYG{n}{max\PYGZus{}power}\PYG{p}{)}

\PYG{n}{boundary\PYGZus{}mode} \PYG{o}{=} \PYG{n+nb}{str}\PYG{p}{(}\PYG{n}{storage}\PYG{o}{.}\PYG{n}{get}\PYG{p}{(}\PYG{l+s+s2}{\PYGZdq{}}\PYG{l+s+s2}{boundary\PYGZus{}mode}\PYG{l+s+s2}{\PYGZdq{}}\PYG{p}{,} \PYG{l+s+s2}{\PYGZdq{}}\PYG{l+s+s2}{cyclic}\PYG{l+s+s2}{\PYGZdq{}}\PYG{p}{)}\PYG{p}{)}
\PYG{n}{loss\PYGZus{}factor} \PYG{o}{=} \PYG{p}{(}\PYG{l+m+mf}{1.0} \PYG{o}{\PYGZhy{}} \PYG{n+nb}{float}\PYG{p}{(}\PYG{n}{storage}\PYG{o}{.}\PYG{n}{get}\PYG{p}{(}\PYG{l+s+s2}{\PYGZdq{}}\PYG{l+s+s2}{self\PYGZus{}loss\PYGZus{}per\PYGZus{}hour}\PYG{l+s+s2}{\PYGZdq{}}\PYG{p}{,} \PYG{l+m+mf}{0.0}\PYG{p}{)}\PYG{p}{)}\PYG{p}{)} \PYG{o}{*}\PYG{o}{*} \PYG{n+nb+bp}{self}\PYG{o}{.}\PYG{n}{dt}
\PYG{n}{eta\PYGZus{}ch} \PYG{o}{=} \PYG{n+nb}{float}\PYG{p}{(}\PYG{n}{storage}\PYG{o}{.}\PYG{n}{get}\PYG{p}{(}\PYG{l+s+s2}{\PYGZdq{}}\PYG{l+s+s2}{charge\PYGZus{}efficiency}\PYG{l+s+s2}{\PYGZdq{}}\PYG{p}{,} \PYG{l+m+mf}{0.95}\PYG{p}{)}\PYG{p}{)}
\PYG{n}{eta\PYGZus{}dis} \PYG{o}{=} \PYG{n+nb}{float}\PYG{p}{(}\PYG{n}{storage}\PYG{o}{.}\PYG{n}{get}\PYG{p}{(}\PYG{l+s+s2}{\PYGZdq{}}\PYG{l+s+s2}{discharge\PYGZus{}efficiency}\PYG{l+s+s2}{\PYGZdq{}}\PYG{p}{,} \PYG{l+m+mf}{0.95}\PYG{p}{)}\PYG{p}{)}
\PYG{k}{if} \PYG{n}{boundary\PYGZus{}mode} \PYG{o+ow}{not} \PYG{o+ow}{in} \PYG{p}{\PYGZob{}}\PYG{l+s+s2}{\PYGZdq{}}\PYG{l+s+s2}{cyclic}\PYG{l+s+s2}{\PYGZdq{}}\PYG{p}{,} \PYG{l+s+s2}{\PYGZdq{}}\PYG{l+s+s2}{fixed}\PYG{l+s+s2}{\PYGZdq{}}\PYG{p}{\PYGZcb{}}\PYG{p}{:}
    \PYG{k}{raise} \PYG{n+ne}{ValueError}\PYG{p}{(}\PYG{l+s+sa}{f}\PYG{l+s+s2}{\PYGZdq{}}\PYG{l+s+s2}{\CJKCode{储热} }\PYG{l+s+si}{\PYGZob{}}\PYG{n}{st\PYGZus{}name}\PYG{l+s+si}{\PYGZcb{}}\PYG{l+s+s2}{ \CJKCode{的} boundary\PYGZus{}mode \CJKCode{仅支持} cyclic \CJKCode{或} fixed}\PYG{l+s+s2}{\PYGZdq{}}\PYG{p}{)}
\PYG{k}{for} \PYG{n}{t} \PYG{o+ow}{in} \PYG{n+nb+bp}{self}\PYG{o}{.}\PYG{n}{times}\PYG{p}{:}
    \PYG{n}{row} \PYG{o}{=} \PYG{p}{\PYGZob{}}
        \PYG{n+nb+bp}{self}\PYG{o}{.}\PYG{n}{variables}\PYG{p}{[}\PYG{l+s+s2}{\PYGZdq{}}\PYG{l+s+s2}{storage\PYGZus{}soc}\PYG{l+s+s2}{\PYGZdq{}}\PYG{p}{]}\PYG{p}{[}\PYG{p}{(}\PYG{n}{st\PYGZus{}name}\PYG{p}{,} \PYG{n}{t}\PYG{p}{)}\PYG{p}{]}\PYG{p}{:} \PYG{l+m+mf}{1.0}\PYG{p}{,}
        \PYG{n+nb+bp}{self}\PYG{o}{.}\PYG{n}{variables}\PYG{p}{[}\PYG{l+s+s2}{\PYGZdq{}}\PYG{l+s+s2}{storage\PYGZus{}charge}\PYG{l+s+s2}{\PYGZdq{}}\PYG{p}{]}\PYG{p}{[}\PYG{p}{(}\PYG{n}{st\PYGZus{}name}\PYG{p}{,} \PYG{n}{t}\PYG{p}{)}\PYG{p}{]}\PYG{p}{:} \PYG{o}{\PYGZhy{}}\PYG{n}{eta\PYGZus{}ch} \PYG{o}{*} \PYG{n+nb+bp}{self}\PYG{o}{.}\PYG{n}{dt}\PYG{p}{,}
        \PYG{n+nb+bp}{self}\PYG{o}{.}\PYG{n}{variables}\PYG{p}{[}\PYG{l+s+s2}{\PYGZdq{}}\PYG{l+s+s2}{storage\PYGZus{}discharge}\PYG{l+s+s2}{\PYGZdq{}}\PYG{p}{]}\PYG{p}{[}\PYG{p}{(}\PYG{n}{st\PYGZus{}name}\PYG{p}{,} \PYG{n}{t}\PYG{p}{)}\PYG{p}{]}\PYG{p}{:} \PYG{n+nb+bp}{self}\PYG{o}{.}\PYG{n}{dt} \PYG{o}{/} \PYG{n}{eta\PYGZus{}dis}\PYG{p}{,}
    \PYG{p}{\PYGZcb{}}
    \PYG{n}{group\PYGZus{}times} \PYG{o}{=} \PYG{n+nb+bp}{self}\PYG{o}{.}\PYG{n}{boundary\PYGZus{}groups}\PYG{p}{[}\PYG{n+nb+bp}{self}\PYG{o}{.}\PYG{n}{time\PYGZus{}boundary\PYGZus{}group}\PYG{p}{[}\PYG{n}{t}\PYG{p}{]}\PYG{p}{]}
    \PYG{n}{is\PYGZus{}group\PYGZus{}start} \PYG{o}{=} \PYG{n}{t} \PYG{o}{==} \PYG{n}{group\PYGZus{}times}\PYG{p}{[}\PYG{l+m+mi}{0}\PYG{p}{]}
    \PYG{k}{if} \PYG{n}{boundary\PYGZus{}mode} \PYG{o}{==} \PYG{l+s+s2}{\PYGZdq{}}\PYG{l+s+s2}{cyclic}\PYG{l+s+s2}{\PYGZdq{}} \PYG{o+ow}{or} \PYG{o+ow}{not} \PYG{n}{is\PYGZus{}group\PYGZus{}start}\PYG{p}{:}
        \PYG{n}{prev} \PYG{o}{=} \PYG{n+nb+bp}{self}\PYG{o}{.}\PYG{n}{time\PYGZus{}predecessor}\PYG{p}{[}\PYG{n}{t}\PYG{p}{]}
        \PYG{n+nb+bp}{self}\PYG{o}{.}\PYG{n}{\PYGZus{}add}\PYG{p}{(}\PYG{n}{row}\PYG{p}{,} \PYG{n+nb+bp}{self}\PYG{o}{.}\PYG{n}{variables}\PYG{p}{[}\PYG{l+s+s2}{\PYGZdq{}}\PYG{l+s+s2}{storage\PYGZus{}soc}\PYG{l+s+s2}{\PYGZdq{}}\PYG{p}{]}\PYG{p}{[}\PYG{p}{(}\PYG{n}{st\PYGZus{}name}\PYG{p}{,} \PYG{n}{prev}\PYG{p}{)}\PYG{p}{]}\PYG{p}{,} \PYG{o}{\PYGZhy{}}\PYG{n}{loss\PYGZus{}factor}\PYG{p}{)}
        \PYG{n}{p}\PYG{o}{.}\PYG{n}{add\PYGZus{}constraint}\PYG{p}{(}\PYG{n}{row}\PYG{p}{,} \PYG{n}{lb}\PYG{o}{=}\PYG{l+m+mf}{0.0}\PYG{p}{,} \PYG{n}{ub}\PYG{o}{=}\PYG{l+m+mf}{0.0}\PYG{p}{)}
    \PYG{k}{else}\PYG{p}{:}
        \PYG{n}{initial\PYGZus{}fraction} \PYG{o}{=} \PYG{n+nb}{float}\PYG{p}{(}\PYG{n}{storage}\PYG{o}{.}\PYG{n}{get}\PYG{p}{(}\PYG{l+s+s2}{\PYGZdq{}}\PYG{l+s+s2}{initial\PYGZus{}soc\PYGZus{}fraction}\PYG{l+s+s2}{\PYGZdq{}}\PYG{p}{,} \PYG{l+m+mf}{0.5}\PYG{p}{)}\PYG{p}{)}
        \PYG{k}{if} \PYG{n}{cap\PYGZus{}idx} \PYG{o+ow}{is} \PYG{o+ow}{not} \PYG{k+kc}{None}\PYG{p}{:}
            \PYG{n}{row}\PYG{p}{[}\PYG{n}{cap\PYGZus{}idx}\PYG{p}{]} \PYG{o}{=} \PYG{o}{\PYGZhy{}}\PYG{n}{loss\PYGZus{}factor} \PYG{o}{*} \PYG{n}{initial\PYGZus{}fraction}
        \PYG{n}{initial\PYGZus{}energy} \PYG{o}{=} \PYG{n}{loss\PYGZus{}factor} \PYG{o}{*} \PYG{n}{initial\PYGZus{}fraction} \PYG{o}{*} \PYG{n}{existing}
        \PYG{n}{p}\PYG{o}{.}\PYG{n}{add\PYGZus{}constraint}\PYG{p}{(}\PYG{n}{row}\PYG{p}{,} \PYG{n}{lb}\PYG{o}{=}\PYG{n}{initial\PYGZus{}energy}\PYG{p}{,} \PYG{n}{ub}\PYG{o}{=}\PYG{n}{initial\PYGZus{}energy}\PYG{p}{)}
\PYG{k}{if} \PYG{n}{boundary\PYGZus{}mode} \PYG{o}{==} \PYG{l+s+s2}{\PYGZdq{}}\PYG{l+s+s2}{fixed}\PYG{l+s+s2}{\PYGZdq{}}\PYG{p}{:}
    \PYG{n}{final\PYGZus{}fraction} \PYG{o}{=} \PYG{n+nb}{float}\PYG{p}{(}\PYG{n}{storage}\PYG{o}{.}\PYG{n}{get}\PYG{p}{(}\PYG{l+s+s2}{\PYGZdq{}}\PYG{l+s+s2}{final\PYGZus{}soc\PYGZus{}min\PYGZus{}fraction}\PYG{l+s+s2}{\PYGZdq{}}\PYG{p}{,} \PYG{l+m+mf}{0.0}\PYG{p}{)}\PYG{p}{)}
    \PYG{k}{for} \PYG{n}{group\PYGZus{}times} \PYG{o+ow}{in} \PYG{n+nb+bp}{self}\PYG{o}{.}\PYG{n}{boundary\PYGZus{}groups}\PYG{o}{.}\PYG{n}{values}\PYG{p}{(}\PYG{p}{)}\PYG{p}{:}
        \PYG{n}{final\PYGZus{}row} \PYG{o}{=} \PYG{p}{\PYGZob{}}\PYG{n+nb+bp}{self}\PYG{o}{.}\PYG{n}{variables}\PYG{p}{[}\PYG{l+s+s2}{\PYGZdq{}}\PYG{l+s+s2}{storage\PYGZus{}soc}\PYG{l+s+s2}{\PYGZdq{}}\PYG{p}{]}\PYG{p}{[}\PYG{p}{(}\PYG{n}{st\PYGZus{}name}\PYG{p}{,} \PYG{n}{group\PYGZus{}times}\PYG{p}{[}\PYG{o}{\PYGZhy{}}\PYG{l+m+mi}{1}\PYG{p}{]}\PYG{p}{)}\PYG{p}{]}\PYG{p}{:} \PYG{l+m+mf}{1.0}\PYG{p}{\PYGZcb{}}
        \PYG{k}{if} \PYG{n}{cap\PYGZus{}idx} \PYG{o+ow}{is} \PYG{o+ow}{not} \PYG{k+kc}{None}\PYG{p}{:}
            \PYG{n}{final\PYGZus{}row}\PYG{p}{[}\PYG{n}{cap\PYGZus{}idx}\PYG{p}{]} \PYG{o}{=} \PYG{o}{\PYGZhy{}}\PYG{n}{final\PYGZus{}fraction}
        \PYG{n}{p}\PYG{o}{.}\PYG{n}{add\PYGZus{}constraint}\PYG{p}{(}\PYG{n}{final\PYGZus{}row}\PYG{p}{,} \PYG{n}{lb}\PYG{o}{=}\PYG{n}{final\PYGZus{}fraction} \PYG{o}{*} \PYG{n}{existing}\PYG{p}{)}
\end{Verbatim}

\endgroup

\Needspace{14\baselineskip}
\subsection*{5．稀疏模型求解}
将目标、变量边界、整数类型及约束矩阵交给SciPy的MILP接口，返回求解器结果供上层复核。

\begingroup
\setstretch{1.0}
\begin{Verbatim}[commandchars=\\\{\},numbers=left,firstnumber=177,stepnumber=1,fontsize=\fontsize{9.5}{11.5}\selectfont,breaklines=true,breakanywhere=true,frame=single,framesep=3mm,numbersep=5pt,rulecolor=\color{black},xleftmargin=18pt]
\PYG{k}{def}\PYG{+w}{ }\PYG{n+nf}{solve}\PYG{p}{(}\PYG{n+nb+bp}{self}\PYG{p}{,} \PYG{n}{mip\PYGZus{}gap}\PYG{p}{:} \PYG{n+nb}{float} \PYG{o}{=} \PYG{l+m+mf}{0.005}\PYG{p}{,} \PYG{n}{time\PYGZus{}limit}\PYG{p}{:} \PYG{n+nb}{float} \PYG{o}{|} \PYG{k+kc}{None} \PYG{o}{=} \PYG{l+m+mf}{120.0}\PYG{p}{)}\PYG{p}{:}
    \PYG{n}{matrix} \PYG{o}{=} \PYG{n+nb+bp}{self}\PYG{o}{.}\PYG{n}{constraint\PYGZus{}matrix}\PYG{p}{(}\PYG{p}{)}
    \PYG{k}{if} \PYG{n+nb+bp}{self}\PYG{o}{.}\PYG{n}{\PYGZus{}compiled\PYGZus{}constraints} \PYG{o+ow}{is} \PYG{k+kc}{None}\PYG{p}{:}
        \PYG{n+nb+bp}{self}\PYG{o}{.}\PYG{n}{\PYGZus{}compiled\PYGZus{}constraints} \PYG{o}{=} \PYG{n}{LinearConstraint}\PYG{p}{(}
            \PYG{n}{matrix}\PYG{p}{,}
            \PYG{n}{np}\PYG{o}{.}\PYG{n}{asarray}\PYG{p}{(}\PYG{n+nb+bp}{self}\PYG{o}{.}\PYG{n}{row\PYGZus{}lower}\PYG{p}{)}\PYG{p}{,}
            \PYG{n}{np}\PYG{o}{.}\PYG{n}{asarray}\PYG{p}{(}\PYG{n+nb+bp}{self}\PYG{o}{.}\PYG{n}{row\PYGZus{}upper}\PYG{p}{)}\PYG{p}{,}
        \PYG{p}{)}
    \PYG{k}{if} \PYG{n+nb+bp}{self}\PYG{o}{.}\PYG{n}{\PYGZus{}compiled\PYGZus{}bounds} \PYG{o+ow}{is} \PYG{k+kc}{None}\PYG{p}{:}
        \PYG{n+nb+bp}{self}\PYG{o}{.}\PYG{n}{\PYGZus{}compiled\PYGZus{}bounds} \PYG{o}{=} \PYG{n}{Bounds}\PYG{p}{(}
            \PYG{n}{np}\PYG{o}{.}\PYG{n}{asarray}\PYG{p}{(}\PYG{n+nb+bp}{self}\PYG{o}{.}\PYG{n}{lower}\PYG{p}{)}\PYG{p}{,} \PYG{n}{np}\PYG{o}{.}\PYG{n}{asarray}\PYG{p}{(}\PYG{n+nb+bp}{self}\PYG{o}{.}\PYG{n}{upper}\PYG{p}{)}
        \PYG{p}{)}
    \PYG{k}{if} \PYG{n+nb+bp}{self}\PYG{o}{.}\PYG{n}{\PYGZus{}compiled\PYGZus{}integrality} \PYG{o+ow}{is} \PYG{k+kc}{None}\PYG{p}{:}
        \PYG{n+nb+bp}{self}\PYG{o}{.}\PYG{n}{\PYGZus{}compiled\PYGZus{}integrality} \PYG{o}{=} \PYG{n}{np}\PYG{o}{.}\PYG{n}{asarray}\PYG{p}{(}\PYG{n+nb+bp}{self}\PYG{o}{.}\PYG{n}{integrality}\PYG{p}{)}
    \PYG{n}{options}\PYG{p}{:} \PYG{n+nb}{dict}\PYG{p}{[}\PYG{n+nb}{str}\PYG{p}{,} \PYG{n+nb}{float} \PYG{o}{|} \PYG{n+nb}{bool}\PYG{p}{]} \PYG{o}{=} \PYG{p}{\PYGZob{}}\PYG{l+s+s2}{\PYGZdq{}}\PYG{l+s+s2}{mip\PYGZus{}rel\PYGZus{}gap}\PYG{l+s+s2}{\PYGZdq{}}\PYG{p}{:} \PYG{n}{mip\PYGZus{}gap}\PYG{p}{,} \PYG{l+s+s2}{\PYGZdq{}}\PYG{l+s+s2}{disp}\PYG{l+s+s2}{\PYGZdq{}}\PYG{p}{:} \PYG{k+kc}{False}\PYG{p}{\PYGZcb{}}
    \PYG{k}{if} \PYG{n}{time\PYGZus{}limit} \PYG{o+ow}{is} \PYG{o+ow}{not} \PYG{k+kc}{None}\PYG{p}{:}
        \PYG{n}{options}\PYG{p}{[}\PYG{l+s+s2}{\PYGZdq{}}\PYG{l+s+s2}{time\PYGZus{}limit}\PYG{l+s+s2}{\PYGZdq{}}\PYG{p}{]} \PYG{o}{=} \PYG{n+nb}{float}\PYG{p}{(}\PYG{n}{time\PYGZus{}limit}\PYG{p}{)}
    \PYG{n+nb+bp}{self}\PYG{o}{.}\PYG{n}{\PYGZus{}last\PYGZus{}solve\PYGZus{}objective} \PYG{o}{=} \PYG{n}{np}\PYG{o}{.}\PYG{n}{asarray}\PYG{p}{(}\PYG{n+nb+bp}{self}\PYG{o}{.}\PYG{n}{objective}\PYG{p}{,} \PYG{n}{dtype}\PYG{o}{=}\PYG{n+nb}{float}\PYG{p}{)}\PYG{o}{.}\PYG{n}{copy}\PYG{p}{(}\PYG{p}{)}
    \PYG{n+nb+bp}{self}\PYG{o}{.}\PYG{n}{\PYGZus{}last\PYGZus{}solve\PYGZus{}fun} \PYG{o}{=} \PYG{k+kc}{None}
    \PYG{n}{result} \PYG{o}{=} \PYG{n}{milp}\PYG{p}{(}
        \PYG{n}{c}\PYG{o}{=}\PYG{n+nb+bp}{self}\PYG{o}{.}\PYG{n}{\PYGZus{}last\PYGZus{}solve\PYGZus{}objective}\PYG{p}{,}
        \PYG{n}{integrality}\PYG{o}{=}\PYG{n+nb+bp}{self}\PYG{o}{.}\PYG{n}{\PYGZus{}compiled\PYGZus{}integrality}\PYG{p}{,}
        \PYG{n}{bounds}\PYG{o}{=}\PYG{n+nb+bp}{self}\PYG{o}{.}\PYG{n}{\PYGZus{}compiled\PYGZus{}bounds}\PYG{p}{,}
        \PYG{n}{constraints}\PYG{o}{=}\PYG{n+nb+bp}{self}\PYG{o}{.}\PYG{n}{\PYGZus{}compiled\PYGZus{}constraints}\PYG{p}{,}
        \PYG{n}{options}\PYG{o}{=}\PYG{n}{options}\PYG{p}{,}
    \PYG{p}{)}
    \PYG{k}{if} \PYG{n}{result}\PYG{o}{.}\PYG{n}{fun} \PYG{o+ow}{is} \PYG{o+ow}{not} \PYG{k+kc}{None}\PYG{p}{:}
        \PYG{n+nb+bp}{self}\PYG{o}{.}\PYG{n}{\PYGZus{}last\PYGZus{}solve\PYGZus{}fun} \PYG{o}{=} \PYG{n+nb}{float}\PYG{p}{(}\PYG{n}{result}\PYG{o}{.}\PYG{n}{fun}\PYG{p}{)}
    \PYG{k}{return} \PYG{n}{result}
\end{Verbatim}

\endgroup

\Needspace{14\baselineskip}
\subsection*{6．固定投资转换}
示出热源和储热的冻结操作。完整函数还处理管网、建筑改造及其他投资项，固定投资复算调用完整函数。

\begingroup
\setstretch{1.0}
\begin{Verbatim}[commandchars=\\\{\},numbers=left,firstnumber=43,stepnumber=1,fontsize=\fontsize{9.5}{11.5}\selectfont,breaklines=true,breakanywhere=true,frame=single,framesep=3mm,numbersep=5pt,rulecolor=\color{black},xleftmargin=18pt]
\PYG{k}{for} \PYG{n}{source} \PYG{o+ow}{in} \PYG{n}{fixed}\PYG{p}{[}\PYG{l+s+s2}{\PYGZdq{}}\PYG{l+s+s2}{sources}\PYG{l+s+s2}{\PYGZdq{}}\PYG{p}{]}\PYG{p}{:}
    \PYG{n}{row} \PYG{o}{=} \PYG{n}{nominal}\PYG{o}{.}\PYG{n}{investments}\PYG{p}{[}
        \PYG{p}{(}\PYG{n}{nominal}\PYG{o}{.}\PYG{n}{investments}\PYG{p}{[}\PYG{l+s+s2}{\PYGZdq{}}\PYG{l+s+s2}{category}\PYG{l+s+s2}{\PYGZdq{}}\PYG{p}{]} \PYG{o}{==} \PYG{l+s+s2}{\PYGZdq{}}\PYG{l+s+s2}{source\PYGZus{}capacity}\PYG{l+s+s2}{\PYGZdq{}}\PYG{p}{)}
        \PYG{o}{\PYGZam{}} \PYG{p}{(}\PYG{n}{nominal}\PYG{o}{.}\PYG{n}{investments}\PYG{p}{[}\PYG{l+s+s2}{\PYGZdq{}}\PYG{l+s+s2}{asset}\PYG{l+s+s2}{\PYGZdq{}}\PYG{p}{]} \PYG{o}{==} \PYG{n}{source}\PYG{p}{[}\PYG{l+s+s2}{\PYGZdq{}}\PYG{l+s+s2}{name}\PYG{l+s+s2}{\PYGZdq{}}\PYG{p}{]}\PYG{p}{)}
    \PYG{p}{]}
    \PYG{n}{addition} \PYG{o}{=} \PYG{n+nb}{float}\PYG{p}{(}\PYG{n}{row}\PYG{p}{[}\PYG{l+s+s2}{\PYGZdq{}}\PYG{l+s+s2}{value}\PYG{l+s+s2}{\PYGZdq{}}\PYG{p}{]}\PYG{o}{.}\PYG{n}{iloc}\PYG{p}{[}\PYG{l+m+mi}{0}\PYG{p}{]}\PYG{p}{)} \PYG{k}{if} \PYG{o+ow}{not} \PYG{n}{row}\PYG{o}{.}\PYG{n}{empty} \PYG{k}{else} \PYG{l+m+mf}{0.0}
    \PYG{n}{retired} \PYG{o}{=} \PYG{n}{nominal}\PYG{o}{.}\PYG{n}{investments}\PYG{p}{[}
        \PYG{p}{(}\PYG{n}{nominal}\PYG{o}{.}\PYG{n}{investments}\PYG{p}{[}\PYG{l+s+s2}{\PYGZdq{}}\PYG{l+s+s2}{category}\PYG{l+s+s2}{\PYGZdq{}}\PYG{p}{]} \PYG{o}{==} \PYG{l+s+s2}{\PYGZdq{}}\PYG{l+s+s2}{source\PYGZus{}retirement}\PYG{l+s+s2}{\PYGZdq{}}\PYG{p}{)}
        \PYG{o}{\PYGZam{}} \PYG{p}{(}\PYG{n}{nominal}\PYG{o}{.}\PYG{n}{investments}\PYG{p}{[}\PYG{l+s+s2}{\PYGZdq{}}\PYG{l+s+s2}{asset}\PYG{l+s+s2}{\PYGZdq{}}\PYG{p}{]} \PYG{o}{==} \PYG{n}{source}\PYG{p}{[}\PYG{l+s+s2}{\PYGZdq{}}\PYG{l+s+s2}{name}\PYG{l+s+s2}{\PYGZdq{}}\PYG{p}{]}\PYG{p}{)}
    \PYG{p}{]}
    \PYG{n}{retirement} \PYG{o}{=} \PYG{n+nb}{float}\PYG{p}{(}\PYG{n}{retired}\PYG{p}{[}\PYG{l+s+s2}{\PYGZdq{}}\PYG{l+s+s2}{value}\PYG{l+s+s2}{\PYGZdq{}}\PYG{p}{]}\PYG{o}{.}\PYG{n}{iloc}\PYG{p}{[}\PYG{l+m+mi}{0}\PYG{p}{]}\PYG{p}{)} \PYG{k}{if} \PYG{o+ow}{not} \PYG{n}{retired}\PYG{o}{.}\PYG{n}{empty} \PYG{k}{else} \PYG{l+m+mf}{0.0}
    \PYG{n}{source}\PYG{p}{[}\PYG{l+s+s2}{\PYGZdq{}}\PYG{l+s+s2}{existing\PYGZus{}capacity\PYGZus{}kw}\PYG{l+s+s2}{\PYGZdq{}}\PYG{p}{]} \PYG{o}{=} \PYG{n+nb}{float}\PYG{p}{(}\PYG{n}{source}\PYG{o}{.}\PYG{n}{get}\PYG{p}{(}\PYG{l+s+s2}{\PYGZdq{}}\PYG{l+s+s2}{existing\PYGZus{}capacity\PYGZus{}kw}\PYG{l+s+s2}{\PYGZdq{}}\PYG{p}{,} \PYG{l+m+mf}{0.0}\PYG{p}{)}\PYG{p}{)} \PYG{o}{+} \PYG{n}{addition} \PYG{o}{\PYGZhy{}} \PYG{n}{retirement}
    \PYG{n}{source}\PYG{p}{[}\PYG{l+s+s2}{\PYGZdq{}}\PYG{l+s+s2}{max\PYGZus{}new\PYGZus{}capacity\PYGZus{}kw}\PYG{l+s+s2}{\PYGZdq{}}\PYG{p}{]} \PYG{o}{=} \PYG{l+m+mf}{0.0}
    \PYG{n}{source}\PYG{p}{[}\PYG{l+s+s2}{\PYGZdq{}}\PYG{l+s+s2}{max\PYGZus{}retire\PYGZus{}capacity\PYGZus{}kw}\PYG{l+s+s2}{\PYGZdq{}}\PYG{p}{]} \PYG{o}{=} \PYG{l+m+mf}{0.0}

\PYG{k}{for} \PYG{n}{storage} \PYG{o+ow}{in} \PYG{n}{fixed}\PYG{p}{[}\PYG{l+s+s2}{\PYGZdq{}}\PYG{l+s+s2}{storages}\PYG{l+s+s2}{\PYGZdq{}}\PYG{p}{]}\PYG{p}{:}
    \PYG{n}{row} \PYG{o}{=} \PYG{n}{nominal}\PYG{o}{.}\PYG{n}{investments}\PYG{p}{[}
        \PYG{p}{(}\PYG{n}{nominal}\PYG{o}{.}\PYG{n}{investments}\PYG{p}{[}\PYG{l+s+s2}{\PYGZdq{}}\PYG{l+s+s2}{category}\PYG{l+s+s2}{\PYGZdq{}}\PYG{p}{]} \PYG{o}{==} \PYG{l+s+s2}{\PYGZdq{}}\PYG{l+s+s2}{storage\PYGZus{}capacity}\PYG{l+s+s2}{\PYGZdq{}}\PYG{p}{)}
        \PYG{o}{\PYGZam{}} \PYG{p}{(}\PYG{n}{nominal}\PYG{o}{.}\PYG{n}{investments}\PYG{p}{[}\PYG{l+s+s2}{\PYGZdq{}}\PYG{l+s+s2}{asset}\PYG{l+s+s2}{\PYGZdq{}}\PYG{p}{]} \PYG{o}{==} \PYG{n}{storage}\PYG{p}{[}\PYG{l+s+s2}{\PYGZdq{}}\PYG{l+s+s2}{name}\PYG{l+s+s2}{\PYGZdq{}}\PYG{p}{]}\PYG{p}{)}
    \PYG{p}{]}
    \PYG{n}{addition} \PYG{o}{=} \PYG{n+nb}{float}\PYG{p}{(}\PYG{n}{row}\PYG{p}{[}\PYG{l+s+s2}{\PYGZdq{}}\PYG{l+s+s2}{value}\PYG{l+s+s2}{\PYGZdq{}}\PYG{p}{]}\PYG{o}{.}\PYG{n}{iloc}\PYG{p}{[}\PYG{l+m+mi}{0}\PYG{p}{]}\PYG{p}{)} \PYG{k}{if} \PYG{o+ow}{not} \PYG{n}{row}\PYG{o}{.}\PYG{n}{empty} \PYG{k}{else} \PYG{l+m+mf}{0.0}
    \PYG{n}{storage}\PYG{p}{[}\PYG{l+s+s2}{\PYGZdq{}}\PYG{l+s+s2}{existing\PYGZus{}capacity\PYGZus{}kwh}\PYG{l+s+s2}{\PYGZdq{}}\PYG{p}{]} \PYG{o}{=} \PYG{n+nb}{float}\PYG{p}{(}\PYG{n}{storage}\PYG{o}{.}\PYG{n}{get}\PYG{p}{(}\PYG{l+s+s2}{\PYGZdq{}}\PYG{l+s+s2}{existing\PYGZus{}capacity\PYGZus{}kwh}\PYG{l+s+s2}{\PYGZdq{}}\PYG{p}{,} \PYG{l+m+mf}{0.0}\PYG{p}{)}\PYG{p}{)} \PYG{o}{+} \PYG{n}{addition}
    \PYG{n}{storage}\PYG{p}{[}\PYG{l+s+s2}{\PYGZdq{}}\PYG{l+s+s2}{max\PYGZus{}new\PYGZus{}capacity\PYGZus{}kwh}\PYG{l+s+s2}{\PYGZdq{}}\PYG{p}{]} \PYG{o}{=} \PYG{l+m+mf}{0.0}
\end{Verbatim}

\endgroup

\Needspace{14\baselineskip}
\subsection*{7．连续供暖季复算}
整理声明服务期，生成或读入固定投资方案，执行连续季节复算，并分别汇总数值、储热和服务期检查。

\begingroup
\setstretch{1.0}
\begin{Verbatim}[commandchars=\\\{\},numbers=left,firstnumber=183,stepnumber=1,fontsize=\fontsize{9.5}{11.5}\selectfont,breaklines=true,breakanywhere=true,frame=single,framesep=3mm,numbersep=5pt,rulecolor=\color{black},xleftmargin=18pt]
\PYG{k}{def}\PYG{+w}{ }\PYG{n+nf}{run\PYGZus{}fixed\PYGZus{}investment\PYGZus{}heating\PYGZus{}season\PYGZus{}replay}\PYG{p}{(}
    \PYG{n}{config}\PYG{p}{:} \PYG{n+nb}{dict}\PYG{p}{[}\PYG{n+nb}{str}\PYG{p}{,} \PYG{n}{Any}\PYG{p}{]}\PYG{p}{,} \PYG{n}{representative\PYGZus{}timeseries}\PYG{p}{:} \PYG{n}{pd}\PYG{o}{.}\PYG{n}{DataFrame}\PYG{p}{,}
    \PYG{n}{chronological\PYGZus{}source}\PYG{p}{:} \PYG{n}{pd}\PYG{o}{.}\PYG{n}{DataFrame}\PYG{p}{,} \PYG{n}{scope}\PYG{p}{:} \PYG{n+nb}{dict}\PYG{p}{[}\PYG{n+nb}{str}\PYG{p}{,} \PYG{n}{Any}\PYG{p}{]}\PYG{p}{,} \PYG{o}{*}\PYG{p}{,}
    \PYG{n}{timestamp\PYGZus{}column}\PYG{p}{:} \PYG{n+nb}{str} \PYG{o}{=} \PYG{l+s+s2}{\PYGZdq{}}\PYG{l+s+s2}{timestamp}\PYG{l+s+s2}{\PYGZdq{}}\PYG{p}{,} \PYG{n}{input\PYGZus{}timezone}\PYG{p}{:} \PYG{n+nb}{str} \PYG{o}{|} \PYG{k+kc}{None} \PYG{o}{=} \PYG{k+kc}{None}\PYG{p}{,}
    \PYG{n}{planning\PYGZus{}time\PYGZus{}limit}\PYG{p}{:} \PYG{n+nb}{float} \PYG{o}{=} \PYG{l+m+mf}{180.0}\PYG{p}{,} \PYG{n}{replay\PYGZus{}time\PYGZus{}limit}\PYG{p}{:} \PYG{n+nb}{float} \PYG{o}{=} \PYG{l+m+mf}{1800.0}\PYG{p}{,}
    \PYG{n}{planning\PYGZus{}mip\PYGZus{}gap}\PYG{p}{:} \PYG{n+nb}{float} \PYG{o}{|} \PYG{k+kc}{None} \PYG{o}{=} \PYG{k+kc}{None}\PYG{p}{,} \PYG{n}{replay\PYGZus{}mip\PYGZus{}gap}\PYG{p}{:} \PYG{n+nb}{float} \PYG{o}{|} \PYG{k+kc}{None} \PYG{o}{=} \PYG{k+kc}{None}\PYG{p}{,}
    \PYG{n}{allow\PYGZus{}cyclic\PYGZus{}state\PYGZus{}assumptions}\PYG{p}{:} \PYG{n+nb}{bool} \PYG{o}{=} \PYG{k+kc}{False}\PYG{p}{,}
    \PYG{n}{fixed\PYGZus{}plan}\PYG{p}{:} \PYG{n}{FixedInvestmentPlanArtifacts} \PYG{o}{|} \PYG{k+kc}{None} \PYG{o}{=} \PYG{k+kc}{None}\PYG{p}{,}
\PYG{p}{)} \PYG{o}{\PYGZhy{}}\PYG{o}{\PYGZgt{}} \PYG{n}{HeatingSeasonReplayArtifacts}\PYG{p}{:}
\PYG{+w}{    }\PYG{l+s+sd}{\PYGZdq{}\PYGZdq{}\PYGZdq{}Plan on the supplied planning horizon and validate only the stated season.\PYGZdq{}\PYGZdq{}\PYGZdq{}}
    \PYG{n}{chronology}\PYG{p}{,} \PYG{n}{scope\PYGZus{}audit} \PYG{o}{=} \PYG{n}{prepare\PYGZus{}heating\PYGZus{}season\PYGZus{}timeseries}\PYG{p}{(}
        \PYG{n}{config}\PYG{p}{,} \PYG{n}{chronological\PYGZus{}source}\PYG{p}{,} \PYG{n}{scope}\PYG{p}{,} \PYG{n}{timestamp\PYGZus{}column}\PYG{o}{=}\PYG{n}{timestamp\PYGZus{}column}\PYG{p}{,} \PYG{n}{input\PYGZus{}timezone}\PYG{o}{=}\PYG{n}{input\PYGZus{}timezone}\PYG{p}{,}
    \PYG{p}{)}
    \PYG{c+c1}{\PYGZsh{} Reject unsupported state assumptions before a potentially expensive solve.}
    \PYG{n}{\PYGZus{}season\PYGZus{}config}\PYG{p}{(}\PYG{n}{config}\PYG{p}{,} \PYG{n}{allow\PYGZus{}cyclic\PYGZus{}state\PYGZus{}assumptions}\PYG{p}{)}
    \PYG{n}{plan} \PYG{o}{=} \PYG{n}{fixed\PYGZus{}plan} \PYG{k}{if} \PYG{n}{fixed\PYGZus{}plan} \PYG{o+ow}{is} \PYG{o+ow}{not} \PYG{k+kc}{None} \PYG{k}{else} \PYG{n}{prepare\PYGZus{}fixed\PYGZus{}investment\PYGZus{}plan}\PYG{p}{(}
        \PYG{n}{config}\PYG{p}{,}\PYG{n}{representative\PYGZus{}timeseries}\PYG{p}{,}\PYG{n}{planning\PYGZus{}time\PYGZus{}limit}\PYG{o}{=}\PYG{n}{planning\PYGZus{}time\PYGZus{}limit}\PYG{p}{,}\PYG{n}{planning\PYGZus{}mip\PYGZus{}gap}\PYG{o}{=}\PYG{n}{planning\PYGZus{}mip\PYGZus{}gap}\PYG{p}{)}
    \PYG{n}{fixed\PYGZus{}config}\PYG{p}{,} \PYG{n}{assumptions} \PYG{o}{=} \PYG{n}{\PYGZus{}season\PYGZus{}config}\PYG{p}{(}\PYG{n}{plan}\PYG{o}{.}\PYG{n}{fixed\PYGZus{}config}\PYG{p}{,} \PYG{n}{allow\PYGZus{}cyclic\PYGZus{}state\PYGZus{}assumptions}\PYG{p}{)}
    \PYG{n}{plan} \PYG{o}{=} \PYG{n}{replace}\PYG{p}{(}\PYG{n}{plan}\PYG{p}{,} \PYG{n}{fixed\PYGZus{}config}\PYG{o}{=}\PYG{n}{fixed\PYGZus{}config}\PYG{p}{)}
    \PYG{n}{\PYGZus{}}\PYG{p}{,} \PYG{n}{scenario\PYGZus{}audit} \PYG{o}{=} \PYG{n}{adapt\PYGZus{}config\PYGZus{}to\PYGZus{}chronology\PYGZus{}scenarios}\PYG{p}{(}\PYG{n}{fixed\PYGZus{}config}\PYG{p}{,} \PYG{n}{chronology}\PYG{p}{)}
    \PYG{n}{preflight} \PYG{o}{=} \PYG{n}{AnnualReplayPreflight}\PYG{p}{(}
        \PYG{n}{chronological\PYGZus{}timeseries}\PYG{o}{=}\PYG{n}{chronology}\PYG{p}{,} \PYG{n}{chronology\PYGZus{}audit}\PYG{o}{=}\PYG{n}{scope\PYGZus{}audit}\PYG{p}{,}
        \PYG{n}{representative\PYGZus{}coverage\PYGZus{}audit}\PYG{o}{=}\PYG{n}{pd}\PYG{o}{.}\PYG{n}{DataFrame}\PYG{p}{(}\PYG{n}{columns}\PYG{o}{=}\PYG{p}{[}\PYG{l+s+s2}{\PYGZdq{}}\PYG{l+s+s2}{coverage\PYGZus{}status}\PYG{l+s+s2}{\PYGZdq{}}\PYG{p}{]}\PYG{p}{)}\PYG{p}{,}
        \PYG{n}{representative\PYGZus{}recommendations}\PYG{o}{=}\PYG{n}{pd}\PYG{o}{.}\PYG{n}{DataFrame}\PYG{p}{(}\PYG{p}{)}\PYG{p}{,} \PYG{n}{scenario\PYGZus{}scope\PYGZus{}audit}\PYG{o}{=}\PYG{n}{scenario\PYGZus{}audit}\PYG{p}{,}
    \PYG{p}{)}
    \PYG{n}{result} \PYG{o}{=} \PYG{n}{replay\PYGZus{}prepared\PYGZus{}fixed\PYGZus{}investment\PYGZus{}plan}\PYG{p}{(}
        \PYG{n}{config}\PYG{p}{,} \PYG{n}{representative\PYGZus{}timeseries}\PYG{p}{,} \PYG{n}{chronology}\PYG{p}{,} \PYG{n}{plan}\PYG{p}{,}
        \PYG{n}{timestamp\PYGZus{}column}\PYG{o}{=}\PYG{n}{timestamp\PYGZus{}column}\PYG{p}{,} \PYG{n}{replay\PYGZus{}time\PYGZus{}limit}\PYG{o}{=}\PYG{n}{replay\PYGZus{}time\PYGZus{}limit}\PYG{p}{,}
        \PYG{n}{replay\PYGZus{}mip\PYGZus{}gap}\PYG{o}{=}\PYG{n}{replay\PYGZus{}mip\PYGZus{}gap}\PYG{p}{,} \PYG{n}{require\PYGZus{}full\PYGZus{}year}\PYG{o}{=}\PYG{k+kc}{False}\PYG{p}{,} \PYG{n}{preflight}\PYG{o}{=}\PYG{n}{preflight}\PYG{p}{,}
    \PYG{p}{)}
    \PYG{c+c1}{\PYGZsh{} Annual planning totals and seasonal totals are intentionally not compared.}
    \PYG{n}{result}\PYG{o}{.}\PYG{n}{kpi\PYGZus{}comparison} \PYG{o}{=} \PYG{n}{pd}\PYG{o}{.}\PYG{n}{DataFrame}\PYG{p}{(}\PYG{p}{)}
    \PYG{n}{storage} \PYG{o}{=} \PYG{n}{seasonal\PYGZus{}storage\PYGZus{}postcheck}\PYG{p}{(}\PYG{n}{result}\PYG{o}{.}\PYG{n}{fixed\PYGZus{}config}\PYG{p}{,} \PYG{n}{result}\PYG{o}{.}\PYG{n}{replay}\PYG{o}{.}\PYG{n}{dispatch}\PYG{p}{,} \PYG{n}{expected\PYGZus{}timeseries}\PYG{o}{=}\PYG{n}{chronology}\PYG{p}{)}
    \PYG{n}{common} \PYG{o}{=} \PYG{n}{annual\PYGZus{}replay\PYGZus{}acceptance\PYGZus{}audit}\PYG{p}{(}
        \PYG{n}{result}\PYG{o}{.}\PYG{n}{fixed\PYGZus{}config}\PYG{p}{,} \PYG{n}{chronology\PYGZus{}audit}\PYG{o}{=}\PYG{n}{scope\PYGZus{}audit}\PYG{p}{,}
        \PYG{n}{representative\PYGZus{}coverage\PYGZus{}audit}\PYG{o}{=}\PYG{n}{preflight}\PYG{o}{.}\PYG{n}{representative\PYGZus{}coverage\PYGZus{}audit}\PYG{p}{,}
        \PYG{n}{investment\PYGZus{}freeze\PYGZus{}audit}\PYG{o}{=}\PYG{n}{result}\PYG{o}{.}\PYG{n}{investment\PYGZus{}freeze\PYGZus{}audit}\PYG{p}{,}
        \PYG{n}{fixed\PYGZus{}investment\PYGZus{}readiness\PYGZus{}audit}\PYG{o}{=}\PYG{n}{result}\PYG{o}{.}\PYG{n}{fixed\PYGZus{}investment\PYGZus{}readiness\PYGZus{}audit}\PYG{p}{,}
        \PYG{n}{planning\PYGZus{}summary}\PYG{o}{=}\PYG{n}{result}\PYG{o}{.}\PYG{n}{planning}\PYG{o}{.}\PYG{n}{summary}\PYG{p}{,} \PYG{n}{replay\PYGZus{}summary}\PYG{o}{=}\PYG{n}{result}\PYG{o}{.}\PYG{n}{replay}\PYG{o}{.}\PYG{n}{summary}\PYG{p}{,}
        \PYG{n}{storage\PYGZus{}boundary\PYGZus{}audit}\PYG{o}{=}\PYG{n}{storage}\PYG{p}{,}
    \PYG{p}{)}
    \PYG{n}{gates} \PYG{o}{=} \PYG{p}{[}\PYG{n}{gate} \PYG{k}{for} \PYG{n}{gate} \PYG{o+ow}{in} \PYG{n}{common}\PYG{p}{[}\PYG{l+s+s2}{\PYGZdq{}}\PYG{l+s+s2}{gates}\PYG{l+s+s2}{\PYGZdq{}}\PYG{p}{]} \PYG{k}{if} \PYG{n}{gate}\PYG{p}{[}\PYG{l+s+s2}{\PYGZdq{}}\PYG{l+s+s2}{gate}\PYG{l+s+s2}{\PYGZdq{}}\PYG{p}{]} \PYG{o+ow}{not} \PYG{o+ow}{in} \PYG{p}{\PYGZob{}}\PYG{l+s+s2}{\PYGZdq{}}\PYG{l+s+s2}{continuous\PYGZus{}full\PYGZus{}calendar\PYGZus{}year}\PYG{l+s+s2}{\PYGZdq{}}\PYG{p}{,} \PYG{l+s+s2}{\PYGZdq{}}\PYG{l+s+s2}{continuous\PYGZus{}full\PYGZus{}service\PYGZus{}year}\PYG{l+s+s2}{\PYGZdq{}}\PYG{p}{,} \PYG{l+s+s2}{\PYGZdq{}}\PYG{l+s+s2}{representative\PYGZus{}period\PYGZus{}coverage}\PYG{l+s+s2}{\PYGZdq{}}\PYG{p}{\PYGZcb{}}\PYG{p}{]}
    \PYG{k}{for} \PYG{n}{gate} \PYG{o+ow}{in} \PYG{n}{gates}\PYG{p}{:}
        \PYG{n}{gate}\PYG{p}{[}\PYG{l+s+s2}{\PYGZdq{}}\PYG{l+s+s2}{gate}\PYG{l+s+s2}{\PYGZdq{}}\PYG{p}{]} \PYG{o}{=} \PYG{n}{gate}\PYG{p}{[}\PYG{l+s+s2}{\PYGZdq{}}\PYG{l+s+s2}{gate}\PYG{l+s+s2}{\PYGZdq{}}\PYG{p}{]}\PYG{o}{.}\PYG{n}{replace}\PYG{p}{(}\PYG{l+s+s2}{\PYGZdq{}}\PYG{l+s+s2}{annual\PYGZus{}}\PYG{l+s+s2}{\PYGZdq{}}\PYG{p}{,} \PYG{l+s+s2}{\PYGZdq{}}\PYG{l+s+s2}{period\PYGZus{}}\PYG{l+s+s2}{\PYGZdq{}}\PYG{p}{)}
        \PYG{n}{gate}\PYG{p}{[}\PYG{l+s+s2}{\PYGZdq{}}\PYG{l+s+s2}{requirement}\PYG{l+s+s2}{\PYGZdq{}}\PYG{p}{]} \PYG{o}{=} \PYG{n}{gate}\PYG{p}{[}\PYG{l+s+s2}{\PYGZdq{}}\PYG{l+s+s2}{requirement}\PYG{l+s+s2}{\PYGZdq{}}\PYG{p}{]}\PYG{o}{.}\PYG{n}{replace}\PYG{p}{(}\PYG{l+s+s2}{\PYGZdq{}}\PYG{l+s+s2}{\CJKCode{全年}}\PYG{l+s+s2}{\PYGZdq{}}\PYG{p}{,} \PYG{l+s+s2}{\PYGZdq{}}\PYG{l+s+s2}{\CJKCode{本服务期}}\PYG{l+s+s2}{\PYGZdq{}}\PYG{p}{)}
    \PYG{n}{gates}\PYG{o}{.}\PYG{n}{insert}\PYG{p}{(}\PYG{l+m+mi}{0}\PYG{p}{,} \PYG{p}{\PYGZob{}}\PYG{l+s+s2}{\PYGZdq{}}\PYG{l+s+s2}{gate}\PYG{l+s+s2}{\PYGZdq{}}\PYG{p}{:} \PYG{l+s+s2}{\PYGZdq{}}\PYG{l+s+s2}{exact\PYGZus{}service\PYGZus{}period}\PYG{l+s+s2}{\PYGZdq{}}\PYG{p}{,} \PYG{l+s+s2}{\PYGZdq{}}\PYG{l+s+s2}{passed}\PYG{l+s+s2}{\PYGZdq{}}\PYG{p}{:} \PYG{n+nb}{bool}\PYG{p}{(}\PYG{n}{scope\PYGZus{}audit}\PYG{p}{[}\PYG{l+s+s2}{\PYGZdq{}}\PYG{l+s+s2}{exact\PYGZus{}scope\PYGZus{}match}\PYG{l+s+s2}{\PYGZdq{}}\PYG{p}{]}\PYG{p}{)}\PYG{p}{,}
                     \PYG{l+s+s2}{\PYGZdq{}}\PYG{l+s+s2}{observed}\PYG{l+s+s2}{\PYGZdq{}}\PYG{p}{:} \PYG{n}{scope\PYGZus{}audit}\PYG{p}{,} \PYG{l+s+s2}{\PYGZdq{}}\PYG{l+s+s2}{requirement}\PYG{l+s+s2}{\PYGZdq{}}\PYG{p}{:} \PYG{l+s+s2}{\PYGZdq{}}\PYG{l+s+s2}{exact timezone\PYGZhy{}aware half\PYGZhy{}open service window}\PYG{l+s+s2}{\PYGZdq{}}\PYG{p}{\PYGZcb{}}\PYG{p}{)}
    \PYG{n}{failed} \PYG{o}{=} \PYG{p}{[}\PYG{n}{gate}\PYG{p}{[}\PYG{l+s+s2}{\PYGZdq{}}\PYG{l+s+s2}{gate}\PYG{l+s+s2}{\PYGZdq{}}\PYG{p}{]} \PYG{k}{for} \PYG{n}{gate} \PYG{o+ow}{in} \PYG{n}{gates} \PYG{k}{if} \PYG{o+ow}{not} \PYG{n}{gate}\PYG{p}{[}\PYG{l+s+s2}{\PYGZdq{}}\PYG{l+s+s2}{passed}\PYG{l+s+s2}{\PYGZdq{}}\PYG{p}{]}\PYG{p}{]}
    \PYG{n}{certificate} \PYG{o}{=} \PYG{p}{\PYGZob{}}\PYG{l+s+s2}{\PYGZdq{}}\PYG{l+s+s2}{schema}\PYG{l+s+s2}{\PYGZdq{}}\PYG{p}{:} \PYG{l+s+s2}{\PYGZdq{}}\PYG{l+s+s2}{heatplan.heating\PYGZus{}season\PYGZus{}certificate.v1}\PYG{l+s+s2}{\PYGZdq{}}\PYG{p}{,} \PYG{l+s+s2}{\PYGZdq{}}\PYG{l+s+s2}{scope\PYGZus{}type}\PYG{l+s+s2}{\PYGZdq{}}\PYG{p}{:} \PYG{l+s+s2}{\PYGZdq{}}\PYG{l+s+s2}{heating\PYGZus{}season}\PYG{l+s+s2}{\PYGZdq{}}\PYG{p}{,}
                   \PYG{l+s+s2}{\PYGZdq{}}\PYG{l+s+s2}{certification\PYGZus{}pass}\PYG{l+s+s2}{\PYGZdq{}}\PYG{p}{:} \PYG{o+ow}{not} \PYG{n}{failed}\PYG{p}{,} \PYG{l+s+s2}{\PYGZdq{}}\PYG{l+s+s2}{gate\PYGZus{}count}\PYG{l+s+s2}{\PYGZdq{}}\PYG{p}{:} \PYG{n+nb}{len}\PYG{p}{(}\PYG{n}{gates}\PYG{p}{)}\PYG{p}{,}
                   \PYG{l+s+s2}{\PYGZdq{}}\PYG{l+s+s2}{passed\PYGZus{}gate\PYGZus{}count}\PYG{l+s+s2}{\PYGZdq{}}\PYG{p}{:} \PYG{n+nb}{len}\PYG{p}{(}\PYG{n}{gates}\PYG{p}{)} \PYG{o}{\PYGZhy{}} \PYG{n+nb}{len}\PYG{p}{(}\PYG{n}{failed}\PYG{p}{)}\PYG{p}{,} \PYG{l+s+s2}{\PYGZdq{}}\PYG{l+s+s2}{failed\PYGZus{}gate\PYGZus{}count}\PYG{l+s+s2}{\PYGZdq{}}\PYG{p}{:} \PYG{n+nb}{len}\PYG{p}{(}\PYG{n}{failed}\PYG{p}{)}\PYG{p}{,}
                   \PYG{l+s+s2}{\PYGZdq{}}\PYG{l+s+s2}{failed\PYGZus{}gates}\PYG{l+s+s2}{\PYGZdq{}}\PYG{p}{:} \PYG{n}{failed}\PYG{p}{,} \PYG{l+s+s2}{\PYGZdq{}}\PYG{l+s+s2}{gates}\PYG{l+s+s2}{\PYGZdq{}}\PYG{p}{:} \PYG{n}{gates}\PYG{p}{,} \PYG{l+s+s2}{\PYGZdq{}}\PYG{l+s+s2}{assumptions}\PYG{l+s+s2}{\PYGZdq{}}\PYG{p}{:} \PYG{n}{assumptions}\PYG{p}{,}
                   \PYG{l+s+s2}{\PYGZdq{}}\PYG{l+s+s2}{annual\PYGZus{}replay\PYGZus{}certified}\PYG{l+s+s2}{\PYGZdq{}}\PYG{p}{:} \PYG{k+kc}{False}\PYG{p}{\PYGZcb{}}
    \PYG{n}{summary} \PYG{o}{=} \PYG{p}{\PYGZob{}}\PYG{l+s+s2}{\PYGZdq{}}\PYG{l+s+s2}{scope}\PYG{l+s+s2}{\PYGZdq{}}\PYG{p}{:} \PYG{n}{scope\PYGZus{}audit}\PYG{p}{,} \PYG{l+s+s2}{\PYGZdq{}}\PYG{l+s+s2}{heating\PYGZus{}season\PYGZus{}certified}\PYG{l+s+s2}{\PYGZdq{}}\PYG{p}{:} \PYG{o+ow}{not} \PYG{n}{failed}\PYG{p}{,}
               \PYG{l+s+s2}{\PYGZdq{}}\PYG{l+s+s2}{annual\PYGZus{}replay\PYGZus{}certified}\PYG{l+s+s2}{\PYGZdq{}}\PYG{p}{:} \PYG{k+kc}{False}\PYG{p}{,} \PYG{l+s+s2}{\PYGZdq{}}\PYG{l+s+s2}{assumptions}\PYG{l+s+s2}{\PYGZdq{}}\PYG{p}{:} \PYG{n}{assumptions}\PYG{p}{,}
               \PYG{l+s+s2}{\PYGZdq{}}\PYG{l+s+s2}{solver\PYGZus{}status}\PYG{l+s+s2}{\PYGZdq{}}\PYG{p}{:} \PYG{n}{result}\PYG{o}{.}\PYG{n}{replay}\PYG{o}{.}\PYG{n}{summary}\PYG{p}{[}\PYG{l+s+s2}{\PYGZdq{}}\PYG{l+s+s2}{solver\PYGZus{}status}\PYG{l+s+s2}{\PYGZdq{}}\PYG{p}{]}\PYG{p}{,}
               \PYG{l+s+s2}{\PYGZdq{}}\PYG{l+s+s2}{unserved\PYGZus{}rate}\PYG{l+s+s2}{\PYGZdq{}}\PYG{p}{:} \PYG{n}{result}\PYG{o}{.}\PYG{n}{replay}\PYG{o}{.}\PYG{n}{summary}\PYG{p}{[}\PYG{l+s+s2}{\PYGZdq{}}\PYG{l+s+s2}{unserved\PYGZus{}rate}\PYG{l+s+s2}{\PYGZdq{}}\PYG{p}{]}\PYG{p}{,}
               \PYG{l+s+s2}{\PYGZdq{}}\PYG{l+s+s2}{numerical\PYGZus{}audit\PYGZus{}passed}\PYG{l+s+s2}{\PYGZdq{}}\PYG{p}{:} \PYG{n}{result}\PYG{o}{.}\PYG{n}{replay}\PYG{o}{.}\PYG{n}{summary}\PYG{o}{.}\PYG{n}{get}\PYG{p}{(}\PYG{l+s+s2}{\PYGZdq{}}\PYG{l+s+s2}{numerical\PYGZus{}audit\PYGZus{}passed}\PYG{l+s+s2}{\PYGZdq{}}\PYG{p}{,} \PYG{k+kc}{False}\PYG{p}{)}\PYG{p}{\PYGZcb{}}
    \PYG{k}{for} \PYG{n}{key} \PYG{o+ow}{in} \PYG{p}{(}\PYG{l+s+s2}{\PYGZdq{}}\PYG{l+s+s2}{demand\PYGZus{}kwh\PYGZus{}per\PYGZus{}year}\PYG{l+s+s2}{\PYGZdq{}}\PYG{p}{,} \PYG{l+s+s2}{\PYGZdq{}}\PYG{l+s+s2}{effective\PYGZus{}heat\PYGZus{}demand\PYGZus{}kwh\PYGZus{}per\PYGZus{}year}\PYG{l+s+s2}{\PYGZdq{}}\PYG{p}{,} \PYG{l+s+s2}{\PYGZdq{}}\PYG{l+s+s2}{unserved\PYGZus{}kwh\PYGZus{}per\PYGZus{}year}\PYG{l+s+s2}{\PYGZdq{}}\PYG{p}{,}
                \PYG{l+s+s2}{\PYGZdq{}}\PYG{l+s+s2}{emissions\PYGZus{}kg\PYGZus{}per\PYGZus{}year}\PYG{l+s+s2}{\PYGZdq{}}\PYG{p}{,} \PYG{l+s+s2}{\PYGZdq{}}\PYG{l+s+s2}{total\PYGZus{}service\PYGZus{}emissions\PYGZus{}kg\PYGZus{}per\PYGZus{}year}\PYG{l+s+s2}{\PYGZdq{}}\PYG{p}{,} \PYG{l+s+s2}{\PYGZdq{}}\PYG{l+s+s2}{electricity\PYGZus{}kwh\PYGZus{}per\PYGZus{}year}\PYG{l+s+s2}{\PYGZdq{}}\PYG{p}{,}
                \PYG{l+s+s2}{\PYGZdq{}}\PYG{l+s+s2}{heating\PYGZus{}cash\PYGZus{}operating\PYGZus{}cost\PYGZus{}yuan\PYGZus{}per\PYGZus{}year}\PYG{l+s+s2}{\PYGZdq{}}\PYG{p}{,} \PYG{l+s+s2}{\PYGZdq{}}\PYG{l+s+s2}{cooling\PYGZus{}demand\PYGZus{}kwh\PYGZus{}per\PYGZus{}year}\PYG{l+s+s2}{\PYGZdq{}}\PYG{p}{)}\PYG{p}{:}
        \PYG{k}{if} \PYG{n}{key} \PYG{o+ow}{in} \PYG{n}{result}\PYG{o}{.}\PYG{n}{replay}\PYG{o}{.}\PYG{n}{summary}\PYG{p}{:}
            \PYG{n}{summary}\PYG{p}{[}\PYG{n}{key}\PYG{o}{.}\PYG{n}{replace}\PYG{p}{(}\PYG{l+s+s2}{\PYGZdq{}}\PYG{l+s+s2}{\PYGZus{}per\PYGZus{}year}\PYG{l+s+s2}{\PYGZdq{}}\PYG{p}{,} \PYG{l+s+s2}{\PYGZdq{}}\PYG{l+s+s2}{\PYGZus{}over\PYGZus{}period}\PYG{l+s+s2}{\PYGZdq{}}\PYG{p}{)}\PYG{p}{]} \PYG{o}{=} \PYG{n}{result}\PYG{o}{.}\PYG{n}{replay}\PYG{o}{.}\PYG{n}{summary}\PYG{p}{[}\PYG{n}{key}\PYG{p}{]}
    \PYG{k}{return} \PYG{n}{HeatingSeasonReplayArtifacts}\PYG{p}{(}\PYG{n}{result}\PYG{p}{,} \PYG{n}{scope\PYGZus{}audit}\PYG{p}{,} \PYG{n}{assumptions}\PYG{p}{,} \PYG{n}{storage}\PYG{p}{,} \PYG{n}{certificate}\PYG{p}{,} \PYG{n}{summary}\PYG{p}{)}
\end{Verbatim}

\endgroup

\Needspace{14\baselineskip}
\subsection*{8．独立碳帽扫描}
费用与排放端点用于设置碳帽区间；各档重新构模、求解和校核，随后筛选样本内非支配解。其投资自由度由输入配置决定。

\begingroup
\setstretch{1.0}
\begin{Verbatim}[commandchars=\\\{\},numbers=left,firstnumber=351,stepnumber=1,fontsize=\fontsize{9.5}{11.5}\selectfont,breaklines=true,breakanywhere=true,frame=single,framesep=3mm,numbersep=5pt,rulecolor=\color{black},xleftmargin=18pt]
\PYG{k}{def}\PYG{+w}{ }\PYG{n+nf}{epsilon\PYGZus{}constraint\PYGZus{}pareto}\PYG{p}{(}
    \PYG{n}{config\PYGZus{}path}\PYG{p}{:} \PYG{n+nb}{str} \PYG{o}{|} \PYG{n}{Path}\PYG{p}{,}
    \PYG{n}{timeseries\PYGZus{}path}\PYG{p}{:} \PYG{n+nb}{str} \PYG{o}{|} \PYG{n}{Path}\PYG{p}{,}
    \PYG{n}{points}\PYG{p}{:} \PYG{n+nb}{int} \PYG{o}{=} \PYG{l+m+mi}{7}\PYG{p}{,}
\PYG{p}{)} \PYG{o}{\PYGZhy{}}\PYG{o}{\PYGZgt{}} \PYG{n}{pd}\PYG{o}{.}\PYG{n}{DataFrame}\PYG{p}{:}
\PYG{+w}{    }\PYG{l+s+sd}{\PYGZdq{}\PYGZdq{}\PYGZdq{}Sample a nondominated system cost/carbon frontier with explicit caps.}

\PYG{l+s+sd}{    The investment and operating model is unchanged across points. A hard reliability}
\PYG{l+s+sd}{    constraint prevents the emissions objective from creating artificial low\PYGZhy{}carbon}
\PYG{l+s+sd}{    solutions through heat curtailment.}
\PYG{l+s+sd}{    \PYGZdq{}\PYGZdq{}\PYGZdq{}}
    \PYG{k}{if} \PYG{n}{points} \PYG{o}{\PYGZlt{}} \PYG{l+m+mi}{2}\PYG{p}{:}
        \PYG{k}{raise} \PYG{n+ne}{ValueError}\PYG{p}{(}\PYG{l+s+s2}{\PYGZdq{}}\PYG{l+s+s2}{points \CJKCode{至少为} 2}\PYG{l+s+s2}{\PYGZdq{}}\PYG{p}{)}
    \PYG{n}{base\PYGZus{}config}\PYG{p}{,} \PYG{n}{frame} \PYG{o}{=} \PYG{n}{load\PYGZus{}inputs}\PYG{p}{(}\PYG{n}{config\PYGZus{}path}\PYG{p}{,} \PYG{n}{timeseries\PYGZus{}path}\PYG{p}{)}
    \PYG{n}{base\PYGZus{}config} \PYG{o}{=} \PYG{n}{copy}\PYG{o}{.}\PYG{n}{deepcopy}\PYG{p}{(}\PYG{n}{base\PYGZus{}config}\PYG{p}{)}
    \PYG{n}{base\PYGZus{}config}\PYG{p}{[}\PYG{l+s+s2}{\PYGZdq{}}\PYG{l+s+s2}{planning}\PYG{l+s+s2}{\PYGZdq{}}\PYG{p}{]}\PYG{p}{[}\PYG{l+s+s2}{\PYGZdq{}}\PYG{l+s+s2}{carbon\PYGZus{}price\PYGZus{}yuan\PYGZus{}per\PYGZus{}kg}\PYG{l+s+s2}{\PYGZdq{}}\PYG{p}{]} \PYG{o}{=} \PYG{l+m+mf}{0.0}
    \PYG{n}{base\PYGZus{}config}\PYG{p}{[}\PYG{l+s+s2}{\PYGZdq{}}\PYG{l+s+s2}{planning}\PYG{l+s+s2}{\PYGZdq{}}\PYG{p}{]}\PYG{o}{.}\PYG{n}{pop}\PYG{p}{(}\PYG{l+s+s2}{\PYGZdq{}}\PYG{l+s+s2}{emissions\PYGZus{}cap\PYGZus{}kg\PYGZus{}per\PYGZus{}year}\PYG{l+s+s2}{\PYGZdq{}}\PYG{p}{,} \PYG{k+kc}{None}\PYG{p}{)}

    \PYG{n}{cost\PYGZus{}optimum} \PYG{o}{=} \PYG{n}{HeatPlanningModel}\PYG{p}{(}\PYG{n}{copy}\PYG{o}{.}\PYG{n}{deepcopy}\PYG{p}{(}\PYG{n}{base\PYGZus{}config}\PYG{p}{)}\PYG{p}{,} \PYG{n}{frame}\PYG{p}{)}\PYG{o}{.}\PYG{n}{solve\PYGZus{}lexicographic}\PYG{p}{(}\PYG{n}{primary}\PYG{o}{=}\PYG{l+s+s2}{\PYGZdq{}}\PYG{l+s+s2}{total\PYGZus{}cost}\PYG{l+s+s2}{\PYGZdq{}}\PYG{p}{)}\PYG{o}{.}\PYG{n}{summary}
    \PYG{n}{emission\PYGZus{}optimum} \PYG{o}{=} \PYG{n}{HeatPlanningModel}\PYG{p}{(}\PYG{n}{copy}\PYG{o}{.}\PYG{n}{deepcopy}\PYG{p}{(}\PYG{n}{base\PYGZus{}config}\PYG{p}{)}\PYG{p}{,} \PYG{n}{frame}\PYG{p}{)}\PYG{o}{.}\PYG{n}{solve\PYGZus{}lexicographic}\PYG{p}{(}\PYG{n}{primary}\PYG{o}{=}\PYG{l+s+s2}{\PYGZdq{}}\PYG{l+s+s2}{emissions}\PYG{l+s+s2}{\PYGZdq{}}\PYG{p}{)}\PYG{o}{.}\PYG{n}{summary}
    \PYG{n}{emission\PYGZus{}column} \PYG{o}{=} \PYG{l+s+s2}{\PYGZdq{}}\PYG{l+s+s2}{total\PYGZus{}service\PYGZus{}emissions\PYGZus{}kg\PYGZus{}per\PYGZus{}year}\PYG{l+s+s2}{\PYGZdq{}}
    \PYG{n}{min\PYGZus{}emissions} \PYG{o}{=} \PYG{n+nb}{float}\PYG{p}{(}\PYG{n}{emission\PYGZus{}optimum}\PYG{p}{[}\PYG{n}{emission\PYGZus{}column}\PYG{p}{]}\PYG{p}{)}
    \PYG{n}{cost\PYGZus{}optimum\PYGZus{}emissions} \PYG{o}{=} \PYG{n+nb}{float}\PYG{p}{(}\PYG{n}{cost\PYGZus{}optimum}\PYG{p}{[}\PYG{n}{emission\PYGZus{}column}\PYG{p}{]}\PYG{p}{)}
    \PYG{k}{if} \PYG{n}{min\PYGZus{}emissions} \PYG{o}{\PYGZgt{}} \PYG{n}{cost\PYGZus{}optimum\PYGZus{}emissions} \PYG{o}{+} \PYG{l+m+mf}{1e\PYGZhy{}6} \PYG{o}{+} \PYG{l+m+mf}{1e\PYGZhy{}8}\PYG{o}{*}\PYG{n+nb}{abs}\PYG{p}{(}\PYG{n}{cost\PYGZus{}optimum\PYGZus{}emissions}\PYG{p}{)}\PYG{p}{:}
        \PYG{k}{raise} \PYG{n+ne}{RuntimeError}\PYG{p}{(}\PYG{l+s+s2}{\PYGZdq{}}\PYG{l+s+s2}{Carbon endpoint exceeds cost endpoint; improve solve budgets before defining an epsilon range}\PYG{l+s+s2}{\PYGZdq{}}\PYG{p}{)}
    \PYG{n}{caps} \PYG{o}{=} \PYG{n}{np}\PYG{o}{.}\PYG{n}{linspace}\PYG{p}{(}\PYG{n}{min\PYGZus{}emissions}\PYG{p}{,} \PYG{n}{cost\PYGZus{}optimum\PYGZus{}emissions}\PYG{p}{,} \PYG{n}{points}\PYG{p}{)}

    \PYG{n}{rows} \PYG{o}{=} \PYG{p}{[}\PYG{p}{]}
    \PYG{k}{for} \PYG{n}{cap} \PYG{o+ow}{in} \PYG{n}{caps}\PYG{p}{:}
        \PYG{n}{config} \PYG{o}{=} \PYG{n}{copy}\PYG{o}{.}\PYG{n}{deepcopy}\PYG{p}{(}\PYG{n}{base\PYGZus{}config}\PYG{p}{)}
        \PYG{c+c1}{\PYGZsh{} A tiny numerical allowance avoids rejecting the emissions\PYGZhy{}optimal endpoint.}
        \PYG{n}{effective\PYGZus{}cap} \PYG{o}{=} \PYG{n+nb}{float}\PYG{p}{(}\PYG{n}{cap}\PYG{p}{)} \PYG{o}{+} \PYG{l+m+mf}{1e\PYGZhy{}7} \PYG{o}{+} \PYG{n+nb}{abs}\PYG{p}{(}\PYG{n+nb}{float}\PYG{p}{(}\PYG{n}{cap}\PYG{p}{)}\PYG{p}{)}\PYG{o}{*}\PYG{l+m+mf}{1e\PYGZhy{}8}
        \PYG{n}{config}\PYG{p}{[}\PYG{l+s+s2}{\PYGZdq{}}\PYG{l+s+s2}{planning}\PYG{l+s+s2}{\PYGZdq{}}\PYG{p}{]}\PYG{p}{[}\PYG{l+s+s2}{\PYGZdq{}}\PYG{l+s+s2}{emissions\PYGZus{}cap\PYGZus{}kg\PYGZus{}per\PYGZus{}year}\PYG{l+s+s2}{\PYGZdq{}}\PYG{p}{]} \PYG{o}{=} \PYG{n}{effective\PYGZus{}cap}
        \PYG{n}{summary} \PYG{o}{=} \PYG{n}{HeatPlanningModel}\PYG{p}{(}\PYG{n}{config}\PYG{p}{,} \PYG{n}{frame}\PYG{p}{)}\PYG{o}{.}\PYG{n}{solve}\PYG{p}{(}\PYG{n}{objective\PYGZus{}mode}\PYG{o}{=}\PYG{l+s+s2}{\PYGZdq{}}\PYG{l+s+s2}{total\PYGZus{}cost}\PYG{l+s+s2}{\PYGZdq{}}\PYG{p}{)}\PYG{o}{.}\PYG{n}{summary}
        \PYG{k}{if} \PYG{o+ow}{not} \PYG{n}{summary}\PYG{p}{[}\PYG{l+s+s2}{\PYGZdq{}}\PYG{l+s+s2}{numerical\PYGZus{}audit\PYGZus{}passed}\PYG{l+s+s2}{\PYGZdq{}}\PYG{p}{]} \PYG{o+ow}{or} \PYG{o+ow}{not} \PYG{n}{summary}\PYG{p}{[}\PYG{l+s+s2}{\PYGZdq{}}\PYG{l+s+s2}{emissions\PYGZus{}cap\PYGZus{}pass}\PYG{l+s+s2}{\PYGZdq{}}\PYG{p}{]}\PYG{p}{:}
            \PYG{k}{raise} \PYG{n+ne}{RuntimeError}\PYG{p}{(}\PYG{l+s+s2}{\PYGZdq{}}\PYG{l+s+s2}{A sampled frontier incumbent failed independent numerical or carbon\PYGZhy{}cap checks}\PYG{l+s+s2}{\PYGZdq{}}\PYG{p}{)}
        \PYG{n}{rows}\PYG{o}{.}\PYG{n}{append}\PYG{p}{(}\PYG{p}{\PYGZob{}}\PYG{l+s+s2}{\PYGZdq{}}\PYG{l+s+s2}{requested\PYGZus{}emissions\PYGZus{}cap\PYGZus{}kg\PYGZus{}per\PYGZus{}year}\PYG{l+s+s2}{\PYGZdq{}}\PYG{p}{:} \PYG{n+nb}{float}\PYG{p}{(}\PYG{n}{cap}\PYG{p}{)}\PYG{p}{,} \PYG{o}{*}\PYG{o}{*}\PYG{n}{summary}\PYG{p}{\PYGZcb{}}\PYG{p}{)}
    \PYG{n}{sampled} \PYG{o}{=} \PYG{n}{pd}\PYG{o}{.}\PYG{n}{json\PYGZus{}normalize}\PYG{p}{(}\PYG{n}{rows}\PYG{p}{)}
    \PYG{n}{result} \PYG{o}{=} \PYG{n}{nondominated\PYGZus{}cost\PYGZus{}emission\PYGZus{}points}\PYG{p}{(}\PYG{n}{sampled}\PYG{p}{)}
    \PYG{n}{result}\PYG{o}{.}\PYG{n}{attrs}\PYG{p}{[}\PYG{l+s+s2}{\PYGZdq{}}\PYG{l+s+s2}{all\PYGZus{}sampled\PYGZus{}points}\PYG{l+s+s2}{\PYGZdq{}}\PYG{p}{]} \PYG{o}{=} \PYG{n}{sampled}\PYG{o}{.}\PYG{n}{to\PYGZus{}dict}\PYG{p}{(}\PYG{n}{orient}\PYG{o}{=}\PYG{l+s+s2}{\PYGZdq{}}\PYG{l+s+s2}{records}\PYG{l+s+s2}{\PYGZdq{}}\PYG{p}{)}
    \PYG{n}{result}\PYG{o}{.}\PYG{n}{attrs}\PYG{p}{[}\PYG{l+s+s2}{\PYGZdq{}}\PYG{l+s+s2}{cost\PYGZus{}endpoint}\PYG{l+s+s2}{\PYGZdq{}}\PYG{p}{]} \PYG{o}{=} \PYG{n}{cost\PYGZus{}optimum}
    \PYG{n}{result}\PYG{o}{.}\PYG{n}{attrs}\PYG{p}{[}\PYG{l+s+s2}{\PYGZdq{}}\PYG{l+s+s2}{carbon\PYGZus{}endpoint}\PYG{l+s+s2}{\PYGZdq{}}\PYG{p}{]} \PYG{o}{=} \PYG{n}{emission\PYGZus{}optimum}
    \PYG{n}{result}\PYG{p}{[}\PYG{l+s+s2}{\PYGZdq{}}\PYG{l+s+s2}{frontier\PYGZus{}scope}\PYG{l+s+s2}{\PYGZdq{}}\PYG{p}{]} \PYG{o}{=} \PYG{l+s+s2}{\PYGZdq{}}\PYG{l+s+s2}{nondominated\PYGZus{}subset\PYGZus{}of\PYGZus{}sampled\PYGZus{}feasible\PYGZus{}solutions}\PYG{l+s+s2}{\PYGZdq{}}
    \PYG{n}{result}\PYG{p}{[}\PYG{l+s+s2}{\PYGZdq{}}\PYG{l+s+s2}{sampled\PYGZus{}point\PYGZus{}count}\PYG{l+s+s2}{\PYGZdq{}}\PYG{p}{]} \PYG{o}{=} \PYG{n+nb}{len}\PYG{p}{(}\PYG{n}{rows}\PYG{p}{)}
    \PYG{n}{result}\PYG{p}{[}\PYG{l+s+s2}{\PYGZdq{}}\PYG{l+s+s2}{removed\PYGZus{}duplicate\PYGZus{}or\PYGZus{}dominated\PYGZus{}count}\PYG{l+s+s2}{\PYGZdq{}}\PYG{p}{]} \PYG{o}{=} \PYG{n+nb}{len}\PYG{p}{(}\PYG{n}{rows}\PYG{p}{)}\PYG{o}{\PYGZhy{}}\PYG{n+nb}{len}\PYG{p}{(}\PYG{n}{result}\PYG{p}{)}
    \PYG{n}{result}\PYG{p}{[}\PYG{l+s+s2}{\PYGZdq{}}\PYG{l+s+s2}{marginal\PYGZus{}abatement\PYGZus{}cost\PYGZus{}yuan\PYGZus{}per\PYGZus{}kg}\PYG{l+s+s2}{\PYGZdq{}}\PYG{p}{]} \PYG{o}{=} \PYG{n}{np}\PYG{o}{.}\PYG{n}{nan}
    \PYG{k}{for} \PYG{n}{i} \PYG{o+ow}{in} \PYG{n+nb}{range}\PYG{p}{(}\PYG{n+nb}{len}\PYG{p}{(}\PYG{n}{result}\PYG{p}{)} \PYG{o}{\PYGZhy{}} \PYG{l+m+mi}{1}\PYG{p}{)}\PYG{p}{:}
        \PYG{n}{delta\PYGZus{}cost} \PYG{o}{=} \PYG{n}{result}\PYG{o}{.}\PYG{n}{loc}\PYG{p}{[}\PYG{n}{i}\PYG{p}{,} \PYG{l+s+s2}{\PYGZdq{}}\PYG{l+s+s2}{objective\PYGZus{}yuan\PYGZus{}per\PYGZus{}year}\PYG{l+s+s2}{\PYGZdq{}}\PYG{p}{]} \PYG{o}{\PYGZhy{}} \PYG{n}{result}\PYG{o}{.}\PYG{n}{loc}\PYG{p}{[}\PYG{n}{i} \PYG{o}{+} \PYG{l+m+mi}{1}\PYG{p}{,} \PYG{l+s+s2}{\PYGZdq{}}\PYG{l+s+s2}{objective\PYGZus{}yuan\PYGZus{}per\PYGZus{}year}\PYG{l+s+s2}{\PYGZdq{}}\PYG{p}{]}
        \PYG{n}{delta\PYGZus{}emission} \PYG{o}{=} \PYG{n}{result}\PYG{o}{.}\PYG{n}{loc}\PYG{p}{[}\PYG{n}{i} \PYG{o}{+} \PYG{l+m+mi}{1}\PYG{p}{,} \PYG{n}{emission\PYGZus{}column}\PYG{p}{]} \PYG{o}{\PYGZhy{}} \PYG{n}{result}\PYG{o}{.}\PYG{n}{loc}\PYG{p}{[}\PYG{n}{i}\PYG{p}{,} \PYG{n}{emission\PYGZus{}column}\PYG{p}{]}
        \PYG{k}{if} \PYG{n}{delta\PYGZus{}emission} \PYG{o}{\PYGZgt{}} \PYG{l+m+mf}{1e\PYGZhy{}6} \PYG{o}{+} \PYG{l+m+mf}{1e\PYGZhy{}8}\PYG{o}{*}\PYG{n+nb}{abs}\PYG{p}{(}\PYG{n}{result}\PYG{o}{.}\PYG{n}{loc}\PYG{p}{[}\PYG{n}{i} \PYG{o}{+} \PYG{l+m+mi}{1}\PYG{p}{,} \PYG{n}{emission\PYGZus{}column}\PYG{p}{]}\PYG{p}{)}\PYG{p}{:}
            \PYG{n}{result}\PYG{o}{.}\PYG{n}{loc}\PYG{p}{[}\PYG{n}{i}\PYG{p}{,} \PYG{l+s+s2}{\PYGZdq{}}\PYG{l+s+s2}{marginal\PYGZus{}abatement\PYGZus{}cost\PYGZus{}yuan\PYGZus{}per\PYGZus{}kg}\PYG{l+s+s2}{\PYGZdq{}}\PYG{p}{]} \PYG{o}{=} \PYG{n}{delta\PYGZus{}cost} \PYG{o}{/} \PYG{n}{delta\PYGZus{}emission}

    \PYG{n}{cost\PYGZus{}span} \PYG{o}{=} \PYG{n}{result}\PYG{p}{[}\PYG{l+s+s2}{\PYGZdq{}}\PYG{l+s+s2}{objective\PYGZus{}yuan\PYGZus{}per\PYGZus{}year}\PYG{l+s+s2}{\PYGZdq{}}\PYG{p}{]}\PYG{o}{.}\PYG{n}{max}\PYG{p}{(}\PYG{p}{)} \PYG{o}{\PYGZhy{}} \PYG{n}{result}\PYG{p}{[}\PYG{l+s+s2}{\PYGZdq{}}\PYG{l+s+s2}{objective\PYGZus{}yuan\PYGZus{}per\PYGZus{}year}\PYG{l+s+s2}{\PYGZdq{}}\PYG{p}{]}\PYG{o}{.}\PYG{n}{min}\PYG{p}{(}\PYG{p}{)}
    \PYG{n}{emission\PYGZus{}span} \PYG{o}{=} \PYG{n}{result}\PYG{p}{[}\PYG{n}{emission\PYGZus{}column}\PYG{p}{]}\PYG{o}{.}\PYG{n}{max}\PYG{p}{(}\PYG{p}{)} \PYG{o}{\PYGZhy{}} \PYG{n}{result}\PYG{p}{[}\PYG{n}{emission\PYGZus{}column}\PYG{p}{]}\PYG{o}{.}\PYG{n}{min}\PYG{p}{(}\PYG{p}{)}
    \PYG{n}{normalized\PYGZus{}cost} \PYG{o}{=} \PYG{p}{(}\PYG{n}{result}\PYG{p}{[}\PYG{l+s+s2}{\PYGZdq{}}\PYG{l+s+s2}{objective\PYGZus{}yuan\PYGZus{}per\PYGZus{}year}\PYG{l+s+s2}{\PYGZdq{}}\PYG{p}{]} \PYG{o}{\PYGZhy{}} \PYG{n}{result}\PYG{p}{[}\PYG{l+s+s2}{\PYGZdq{}}\PYG{l+s+s2}{objective\PYGZus{}yuan\PYGZus{}per\PYGZus{}year}\PYG{l+s+s2}{\PYGZdq{}}\PYG{p}{]}\PYG{o}{.}\PYG{n}{min}\PYG{p}{(}\PYG{p}{)}\PYG{p}{)} \PYG{o}{/} \PYG{n+nb}{max}\PYG{p}{(}\PYG{n}{cost\PYGZus{}span}\PYG{p}{,} \PYG{l+m+mf}{1e\PYGZhy{}9}\PYG{p}{)}
    \PYG{n}{normalized\PYGZus{}emission} \PYG{o}{=} \PYG{p}{(}\PYG{n}{result}\PYG{p}{[}\PYG{n}{emission\PYGZus{}column}\PYG{p}{]} \PYG{o}{\PYGZhy{}} \PYG{n}{result}\PYG{p}{[}\PYG{n}{emission\PYGZus{}column}\PYG{p}{]}\PYG{o}{.}\PYG{n}{min}\PYG{p}{(}\PYG{p}{)}\PYG{p}{)} \PYG{o}{/} \PYG{n+nb}{max}\PYG{p}{(}\PYG{n}{emission\PYGZus{}span}\PYG{p}{,} \PYG{l+m+mf}{1e\PYGZhy{}9}\PYG{p}{)}
    \PYG{n}{result}\PYG{p}{[}\PYG{l+s+s2}{\PYGZdq{}}\PYG{l+s+s2}{distance\PYGZus{}to\PYGZus{}ideal}\PYG{l+s+s2}{\PYGZdq{}}\PYG{p}{]} \PYG{o}{=} \PYG{n}{np}\PYG{o}{.}\PYG{n}{sqrt}\PYG{p}{(}\PYG{n}{normalized\PYGZus{}cost}\PYG{o}{*}\PYG{o}{*}\PYG{l+m+mi}{2} \PYG{o}{+} \PYG{n}{normalized\PYGZus{}emission}\PYG{o}{*}\PYG{o}{*}\PYG{l+m+mi}{2}\PYG{p}{)}
    \PYG{n}{result}\PYG{p}{[}\PYG{l+s+s2}{\PYGZdq{}}\PYG{l+s+s2}{recommended\PYGZus{}compromise}\PYG{l+s+s2}{\PYGZdq{}}\PYG{p}{]} \PYG{o}{=} \PYG{k+kc}{False}
    \PYG{n}{result}\PYG{o}{.}\PYG{n}{loc}\PYG{p}{[}\PYG{n}{result}\PYG{p}{[}\PYG{l+s+s2}{\PYGZdq{}}\PYG{l+s+s2}{distance\PYGZus{}to\PYGZus{}ideal}\PYG{l+s+s2}{\PYGZdq{}}\PYG{p}{]}\PYG{o}{.}\PYG{n}{idxmin}\PYG{p}{(}\PYG{p}{)}\PYG{p}{,} \PYG{l+s+s2}{\PYGZdq{}}\PYG{l+s+s2}{recommended\PYGZus{}compromise}\PYG{l+s+s2}{\PYGZdq{}}\PYG{p}{]} \PYG{o}{=} \PYG{k+kc}{True}
    \PYG{n}{result}\PYG{p}{[}\PYG{l+s+s2}{\PYGZdq{}}\PYG{l+s+s2}{recommendation\PYGZus{}method}\PYG{l+s+s2}{\PYGZdq{}}\PYG{p}{]} \PYG{o}{=} \PYG{l+s+s2}{\PYGZdq{}}\PYG{l+s+s2}{closest\PYGZus{}normalized\PYGZus{}ideal\PYGZus{}point}\PYG{l+s+s2}{\PYGZdq{}}
    \PYG{k}{return} \PYG{n}{result}
\end{Verbatim}

\endgroup
